% ROZDZIAL 14 — Type erasure, typy sumaryczne i interfejsy generyczne (C++20).
% Fragment do zbiorczego skryptu, BEZ preambuly i \begin{document}.
% Wymagane: listings, amsmath, amssymb, booktabs, tabularx, array, hyperref.
% Dolaczenie: % ROZDZIAL 14 — Type erasure, typy sumaryczne i interfejsy generyczne (C++20).
% Fragment do zbiorczego skryptu, BEZ preambuly i \begin{document}.
% Wymagane: listings, amsmath, amssymb, booktabs, tabularx, array, hyperref.
% Dolaczenie: % ROZDZIAL 14 — Type erasure, typy sumaryczne i interfejsy generyczne (C++20).
% Fragment do zbiorczego skryptu, BEZ preambuly i \begin{document}.
% Wymagane: listings, amsmath, amssymb, booktabs, tabularx, array, hyperref.
% Dolaczenie: % ROZDZIAL 14 — Type erasure, typy sumaryczne i interfejsy generyczne (C++20).
% Fragment do zbiorczego skryptu, BEZ preambuly i \begin{document}.
% Wymagane: listings, amsmath, amssymb, booktabs, tabularx, array, hyperref.
% Dolaczenie: \input{rozdzial_14_type_erasure_typy_sumaryczne_cpp.tex}
% Kompletne programy sa oznaczone w podpisach listingow.

\chapter{Wymazywanie typów i typy sumaryczne}
\label{chap:cpp14-type-erasure}

\section{Problem: jak reprezentować różne rodzaje wartości?}

W programowaniu generycznym często chcemy posługiwać się różnymi typami
przez wspólny interfejs. C++ oferuje kilka mechanizmów realizacji tego celu,
ale mechanizmy te nie są zamienne.
Wyobraźmy sobie aplikację z trzema poleceniami: zapisaniem pliku,
dodaniem dwóch liczb i zakończeniem pracy. Gdy znamy wszystkie trzy
rodzaje poleceń, możemy je wymienić w \texttt{std::variant}. Jeśli
zewnętrzne moduły mają dodawać nowe polecenia bez zmiany tej listy,
potrzebujemy otwartego interfejsu, np. klasy bazowej z funkcją wirtualną.
Wybór mechanizmu zaczynamy więc od pytań:
\begin{enumerate}
\item Czy zestaw dozwolonych typów jest \emph{zamknięty} (ustalony z góry),
      czy \emph{otwarty} (może zostać rozszerzony o nowy typ)?
\item Czy chcemy przechować \emph{wartość}, czy jedynie referencję/uchwyt
      do istniejącego obiektu?
\item Czy wywoływane operacje mają być znane już podczas kompilacji,
      czy wybierane dynamicznie w czasie działania programu?
\item Czy dopuszczamy alokację dynamiczną, kopiowanie oraz koszt
      pośredniego wywołania funkcji?
\end{enumerate}

\emph{Typ sumaryczny} (ang. \emph{sum type}) reprezentuje alternatywę:
wartość należy do jednego z kilku przypadków. W C++17 taką rolę pełni
\texttt{std::variant<Ts...>}; \texttt{std::optional<T>} opisuje wariant
,,brak wartości albo wartość \texttt{T}''. Natomiast \emph{type erasure}
(\emph{wymazywanie typu}) pozwala ukryć konkretny typ obiektu za
interfejsem nieparametryzowanym przez ten typ. Standardowe przykłady
to \texttt{std::any} i \texttt{std::function}.

\begin{center}
\small
\begin{tabularx}{\textwidth}{@{}lXX@{}}
\toprule
\textbf{Mechanizm} & \textbf{Informacja o typie} & \textbf{Zasadnicze zastosowanie} \\
\midrule
\texttt{optional<T>} & Jedna możliwość \texttt{T} lub brak & Wynik opcjonalny, wartość nieobecna \\
\texttt{variant<Ts...>} & Zamknięty zbiór alternatyw & Zdarzenia, komendy, model stanu \\
\texttt{any} & Typ ustalany w czasie wykonania & Heterogeniczne wartości bez wspólnego interfejsu operacji \\
\texttt{function<R(Args...)>} & Dowolny zgodny obiekt wywoływalny & Przechowywany callback, funkcja jako wartość \\
\texttt{unique\_ptr<Base>} & Otwarta hierarchia klas & Polimorfizm dynamiczny i tożsamość obiektu \\
Szablon/\texttt{concept} & Typ znany przy instancjowaniu & Polimorfizm statyczny, specjalizacja kodu \\
\bottomrule
\end{tabularx}
\end{center}

Najprostsza ściąga to: \texttt{optional<int>} oznacza liczbę albo jej
brak, \texttt{variant<int, std::string>} jedną z dwóch znanych
postaci wartości, a \texttt{any} wartość o typie, którego autor
pojemnika nie musi znać. \texttt{function<void()>} przechowuje
różne działania, które można wywołać w taki sam sposób. W kolejnych
sekcjach zobaczymy, co z tych wyborów wynika dla sprawdzania typów,
własności danych i kosztu wykonania.

\section{\texttt{std::optional<T>}: wartość, której może nie być}

\begin{center}\small
\begin{tabularx}{\textwidth}{@{}p{0.33\textwidth}X@{}}
\toprule
Typ lub operacja & Co reprezentuje \\
\midrule
\code{std::optional<T>} & Wartość typu \code{T} albo jej brak; nie opisuje przyczyny błędu. \\
\code{has_value()} & Sprawdza, czy \code{optional} przechowuje wartość. \\
\code{value()} & Zwraca wartość lub zgłasza \code{std::bad_optional_access}, gdy jej brak. \\
\code{value_or(fallback)} & Zwraca wartość albo podany zamiennik, konwertowany do \code{T}. \\
\code{std::variant<Ts...>} & Dokładnie jedną z ustalonych alternatyw typów. \\
\code{std::any} & Wartość typu ustalonego dopiero w czasie wykonania; przy odczycie trzeba sprawdzić typ. \\
\bottomrule
\end{tabularx}
\end{center}
\label{sec:cpp14-optional}

Typ \texttt{std::optional<T>} (C++17, nagłówek \texttt{<optional>})
przechowuje bezpośrednio obiekt typu \texttt{T}, jeżeli jest on obecny,
lub reprezentuje stan pusty. Jest to pełnoprawna \emph{wartość}, a nie
wskaźnik. W szczególności pusty \texttt{optional} nie oznacza wskaźnika
zerowego, a niepusty nie wymaga automatycznie osobnej alokacji dynamicznej.

\begin{lstlisting}[caption={Najważniejsze operacje optional (C++17); kompletny program.}]
#include <cassert>
#include <optional>
#include <string>

int main() {
    std::optional<int> x;         // stan pusty
    assert(!x.has_value());
    assert(x == std::nullopt);

    x = 42;
    assert(x.has_value());
    assert(*x == 42);
    assert(x.value() == 42);
    assert(x.value_or(-1) == 42);

    x.reset();
    assert(x.value_or(-1) == -1);
    x.emplace(7);
    assert(*x == 7);

    std::optional<std::string> name{std::in_place, 3, 'a'};
    assert(*name == "aaa");
}
\end{lstlisting}

Dostęp przez \texttt{*x} i \texttt{x->...} wymaga obecności wartości:
\emph{użycie pustego obiektu w taki sposób jest niezdefiniowanym
zachowaniem}. Metoda \texttt{value()} wykonuje sprawdzenie i dla stanu
pustego zgłasza \texttt{std::bad\_optional\_access}. Warto pamiętać,
że \texttt{value\_or(default)} przyjmuje wartość domyślną jako argument
funkcji: wyrażenie tworzące argument jest obliczane również wtedy,
gdy \texttt{optional} zawiera wartość. Może ponadto następować kopiowanie
lub przeniesienie przechowywanego obiektu.

\subsection{Praktyczny przykład: parsowanie liczby}

Funkcja parsująca dane wejściowe może zwrócić \texttt{optional}, jeśli
interesuje nas tylko informacja ,,udało się/nie udało się''.
\texttt{std::from\_chars} unika wyjątków i nie wymaga tworzenia dodatkowego
łańcucha zakończonego zerem. W przykładzie wymagamy, aby \emph{cały}
łańcuch był prawidłową liczbą.

\begin{lstlisting}[caption={Parsowanie int do optional (C++17); kompletny program.}]
#include <cassert>
#include <charconv>
#include <optional>
#include <string_view>
#include <system_error>

std::optional<int> parse_int(std::string_view text) {
    int result = 0;
    auto [ptr, ec] = std::from_chars(
        text.data(), text.data() + text.size(), result);
    if (ec != std::errc{} || ptr != text.data() + text.size())
        return std::nullopt;
    return result;
}

int main() {
    assert(parse_int("123") == 123);
    assert(parse_int("123x") == std::nullopt);
    assert(parse_int("") == std::nullopt);
    assert(parse_int("99999999999999999999") == std::nullopt);
}
\end{lstlisting}

\subsection{Własność, referencje i koszt}

\texttt{optional<T>} posiada swoją wartość i zwykle zajmuje więcej pamięci
niż sam \texttt{T}, ponieważ musi przechowywać informację, czy wartość
istnieje. Rozmiar zależy od typu i implementacji. Standard C++20
nie udostępnia \texttt{std::optional<T\&>}. Jeśli wynik ma wskazywać na
istniejący obiekt, można użyć wskaźnika nieposiadającego własności albo
\texttt{std::optional<std::reference\_wrapper<T>>}. W obu przypadkach
programista odpowiada za czas życia wskazywanego obiektu.

\begin{lstlisting}[caption={Opcjonalna referencja: brak przenoszenia własności; fragment.}]
#include <functional>
#include <optional>
#include <string>

std::optional<std::reference_wrapper<const std::string>>
find_name(bool exists, const std::string& name) {
    if (exists) return std::cref(name);
    return std::nullopt;
}
// Zwracana referencja jest poprawna tylko gdy `name` nadal istnieje.
\end{lstlisting}

Gdy brak wartości jest sytuacją normalną, \texttt{optional} najczęściej
czytelniej wyraża intencję niż wyjątkowe zakończenie funkcji. Jeżeli
konieczna jest informacja \emph{dlaczego} nie uzyskano wyniku, należy
zaprojektować osobny typ błędu, użyć \texttt{variant} lub (w C++23)
\texttt{std::expected<T,E>}.

\section{\texttt{std::variant<Ts...>}: zamknięta alternatywa typów}
\label{sec:cpp14-variant}

\texttt{std::variant<Ts...>} (C++17, nagłówek \texttt{<variant>})
przechowuje \emph{jedną aktywną alternatywę} z listy typów \texttt{Ts...}.
Nie jest odpowiednikiem \texttt{std::tuple<Ts...>}: krotka zawiera
\emph{wszystkie} składowe, a wariant dokładnie jeden z dopuszczonych
przypadków (poza szczególnym stanem \texttt{valueless\_by\_exception()}).

\begin{lstlisting}[caption={Tworzenie i badanie aktywnej alternatywy (C++17); kompletny program.}]
#include <cassert>
#include <string>
#include <variant>

int main() {
    using Value = std::variant<int, std::string>;
    Value x = 7;
    assert(x.index() == 0);
    assert(std::holds_alternative<int>(x));
    assert(std::get<int>(x) == 7);

    x = std::string("hello");
    assert(x.index() == 1);
    assert(std::get_if<int>(&x) == nullptr);
    assert(std::get_if<std::string>(&x) != nullptr);

    x.emplace<0>(19);
    assert(std::get<0>(x) == 19);
}
\end{lstlisting}

Niepoprawne \texttt{std::get<T>(x)} (gdy aktywny jest inny typ)
zgłasza \texttt{std::bad\_variant\_access}. Alternatywa
\texttt{std::get\_if<T>(\&x)} zwraca \texttt{nullptr}, a zatem nadaje się
do sprawdzania bez wyjątku. \texttt{std::get<T>} i
\texttt{std::holds\_alternative<T>} wymagają, aby \texttt{T} występował
\emph{dokładnie raz} w liście alternatyw; dla powtarzających się typów
używamy indeksów. Nie należy także zapominać o możliwym niepowodzeniu
operacji konwersji przy konstruowaniu wariantu.

\subsection{Domyślna alternatywa i \texttt{std::monostate}}

Domyślnie konstruowany wariant próbuje skonstruować pierwszą alternatywę.
Jeżeli chcemy jawnie reprezentować stan pusty, można umieścić na początku
\texttt{std::monostate}; zapewnia on także domyślną konstruowalność
wariantu, gdy inny pierwszy typ nie jest domyślnie konstruowalny.

\begin{lstlisting}[caption={Jawny stan początkowy monostate; kompletny program.}]
#include <cassert>
#include <string>
#include <variant>

int main() {
    using State = std::variant<std::monostate, int, std::string>;
    State s;
    assert(std::holds_alternative<std::monostate>(s));
    s = std::string("ready");
    assert(std::get<std::string>(s) == "ready");
}
\end{lstlisting}

\subsection{Odwiedzanie wariantu: \texttt{std::visit}}

Samo wielokrotne sprawdzanie \texttt{holds\_alternative} szybko
prowadzi do rozbudowanych instrukcji \texttt{if}. Mechanizm
\texttt{std::visit(visitor, variant)} wywołuje odwiedzającego z aktywną
wartością odpowiedniego typu. Gdy odwiedzający jest lambdą generyczną,
kompilator instancjuje jej operator wywołania dla różnych alternatyw.

\begin{lstlisting}[caption={Jeden odwiedzający generyczny; kompletny program.}]
#include <cassert>
#include <cstddef>
#include <string>
#include <type_traits>
#include <variant>

int main() {
    std::variant<int, double, std::string> v = "hi";

    const std::size_t size = std::visit([](const auto& x) {
        using T = std::decay_t<decltype(x)>;
        if constexpr (std::is_same_v<T, std::string>)
            return x.size();
        else
            return sizeof(x);
    }, v);

    assert(size == 2);
}
\end{lstlisting}

Do osobnej obsługi poszczególnych typów wygodny jest wzorzec
\texttt{overloaded}, znany również z rozdziału o lambdach.
Wykorzystuje on dziedziczenie po kilku obiektach funkcyjnych oraz
ujawnienie wszystkich operatorów wywołania za pomocą \texttt{using}.

\begin{lstlisting}[caption={overloaded oraz std::visit (C++17); kompletny program.}]
#include <cassert>
#include <cstddef>
#include <string>
#include <variant>

template<class... Fs>
struct overloaded : Fs... {
    using Fs::operator()...;
};
template<class... Fs> overloaded(Fs...) -> overloaded<Fs...>;

int main() {
    using Data = std::variant<int, double, std::string>;
    Data data = std::string("abc");

    auto length = std::visit(overloaded{
        [](int) { return std::size_t{1}; },
        [](double) { return std::size_t{1}; },
        [](const std::string& s) { return s.size(); }
    }, data);
    assert(length == 3);
}
\end{lstlisting}

W przypadku odwiedzania zwracającego wartość wymagany jest zgodny typ
wyniku wszystkich możliwych wywołań (w praktyce należy zapewnić
identyczny typ i kategorię wartości zwracanych wyrażeń, z wyjątkiem
szczególnych reguł dla \texttt{void}). Przy wielu wariantach
\texttt{std::visit} uwzględnia \emph{kombinacje} alternatyw. To daje
czytelny kod przy małej liczbie przypadków, ale może zwiększyć liczbę
instancjacji szablonów.

\begin{lstlisting}[caption={Odwiedzenie dwóch wariantów naraz; kompletny program.}]
#include <cassert>
#include <variant>

int main() {
    std::variant<int, double> a = 3;
    std::variant<int, double> b = 2.5;
    const double sum = std::visit(
        [](const auto& x, const auto& y) -> double {
            return static_cast<double>(x) + static_cast<double>(y);
        }, a, b);
    assert(sum == 5.5);
}
\end{lstlisting}

\subsection{Modyfikacja stanu i kwalifikacja \texttt{const}}

\texttt{std::visit} respektuje typ i kategorię wartości przekazanego
wariantu. Odwiedzając obiekt przez \texttt{variant\&}, możemy modyfikować
aktywną alternatywę; przez \texttt{const variant\&} uzyskujemy
referencję stałą. Przy przekazaniu \texttt{std::move(v)} odwiedzający
może otrzymać referencję rvalue i przenieść przechowywaną wartość.
Przeniesienie nie musi zmieniać aktywnego indeksu wariantu, ale zawartość
przeniesionego obiektu może zostać zmieniona zgodnie z semantyką typu.

\begin{lstlisting}[caption={Modyfikacja aktywnej alternatywy; kompletny program.}]
#include <cassert>
#include <variant>

int main() {
    std::variant<int, double> v = 10;
    std::visit([](auto& x) { x += 1; }, v);
    assert(std::get<int>(v) == 11);
}
\end{lstlisting}

\subsection{Stan \texttt{valueless\_by\_exception}}

Wyjątek podczas przypisania lub tworzenia nowej wartości może
pozostawić \texttt{variant} bez aktywnej alternatywy. Ten stan
wykrywa metoda \texttt{valueless\_by\_exception()}.
Wtedy \texttt{index()} zwraca \texttt{std::variant\_npos}.
Wywołanie \texttt{std::visit} zgłasza \texttt{std::bad\_variant\_access}.
Nie jest to zwykła ,,pusta'' alternatywa. Jeśli jest potrzebna,
można użyć \texttt{std::monostate}.

\paragraph{Ważne ograniczenia.}
\texttt{variant} nie przechowuje typów referencyjnych ani \texttt{void};
nie można także dodać dowolnego nowego typu bez zmiany listy alternatyw.
Duży wariant może zajmować pamięć odpowiadającą największej alternatywie
plus informację o aktywnym przypadku i ewentualne wyrównanie; dokładny
rozmiar nie jest stałą wielkością określoną przez standard.

\section{Przykład projektowy: interpreter prostych poleceń}
\label{sec:cpp14-commands}

Rozpatrzmy miniaturowy model poleceń. Każdy rodzaj polecenia ma inne pola,
ale wszystkie należą do zamkniętej rodziny. \texttt{variant} oddziela
\emph{dane} od \emph{operacji}, a \texttt{visit} pozwala jednoznacznie
obsłużyć każdy przypadek.

\begin{lstlisting}[caption={Modelowanie poleceń za pomocą variant; kompletny program.}]
#include <cassert>
#include <string>
#include <variant>
#include <vector>

struct Add  { int a, b; };
struct Save { std::string filename; };
struct Quit {};
using Command = std::variant<Add, Save, Quit>;

struct Execute {
    int& total;
    bool& stopped;

    void operator()(const Add& c) const {
        total += c.a + c.b;
    }
    void operator()(const Save& c) const {
        if (c.filename.empty()) stopped = true;
    }
    void operator()(const Quit&) const {
        stopped = true;
    }
};

int main() {
    std::vector<Command> commands{
        Add{2, 3}, Save{"result.txt"}, Add{4, 5}, Quit{}
    };
    int total = 0;
    bool stopped = false;
    for (const auto& command : commands) {
        if (stopped) break;
        std::visit(Execute{total, stopped}, command);
    }
    assert(total == 14);
    assert(stopped);
}
\end{lstlisting}

Dodanie nowego typu polecenia, np. \texttt{Delete}, wymaga rozszerzenia
\texttt{Command} i obsługi nowego przypadku. W dobrze skonstruowanym
odwiedzającym brakująca obsługa będzie błędem kompilacji. To zaleta
typów sumarycznych przy modelowaniu \emph{zamkniętego} katalogu zdarzeń.

\section{\texttt{std::any}: wartość o nieznanym typie}
\label{sec:cpp14-any}

\texttt{std::any} (C++17, nagłówek \texttt{<any>}) może przechowywać
obiekt dowolnego typu spełniającego wymaganie kopiowalności i pozwala
później odzyskać dokładny typ przez \texttt{std::any\_cast}. To klasyczna
implementacja \emph{type erasure}: użytkownik obiektu \texttt{any}
nie zna szablonowego typu przechowywanej wartości.

\begin{lstlisting}[caption={std::any i bezpieczne wydobycie typu (C++17); kompletny program.}]
#include <any>
#include <cassert>
#include <string>

int main() {
    std::any value = 17;
    assert(value.has_value());
    assert(std::any_cast<int>(value) == 17);

    if (auto ptr = std::any_cast<int>(&value)) {
        *ptr += 1;
    }
    assert(std::any_cast<int>(value) == 18);

    value.emplace<std::string>("text");
    assert(std::any_cast<int>(&value) == nullptr);
    assert(std::any_cast<const std::string&>(value) == "text");

    value.reset();
    assert(!value.has_value());
}
\end{lstlisting}

\texttt{std::any\_cast<T>(value)} zgłasza \texttt{std::bad\_any\_cast},
jeśli aktywny typ nie odpowiada żądanemu. Wariant wskaźnikowy
\texttt{std::any\_cast<T>(\&value)} zwraca \texttt{nullptr}; jest zatem
wygodny do ostrożnego sprawdzania typu. \texttt{type()} udostępnia
\texttt{std::type\_info} opisujący aktualnie przechowywany typ,
a dla pustego obiektu zwraca \texttt{typeid(void)}.

\begin{lstlisting}[caption={\texttt{any\_cast} nie jest konwersją numeryczną; kompletny program.}]
#include <any>
#include <cassert>
#include <typeinfo>

int main() {
    std::any value = 42;
    assert(value.type() == typeid(int));
    assert(std::any_cast<double>(&value) == nullptr);
    assert(std::any_cast<int>(&value) != nullptr);
}
\end{lstlisting}

\texttt{any} przechowuje \emph{wartość}, nie udostępnia jednak uniwersalnego
operatora \texttt{+}, porównania czy drukowania tej wartości. Aby wykonać
operację, trzeba znać typ albo dostarczyć odrębny mechanizm dyspozycji.
Ponadto przechowywany typ musi być kopiowalny: \texttt{std::any}
nie nadaje się do bezpośredniego przechowywania \texttt{std::unique\_ptr<T>}.
Niektóre implementacje wykorzystują \emph{small-object optimization},
ale standard nie obiecuje uniknięcia alokacji dla konkretnego małego typu.

\paragraph{Kiedy używać?}
\texttt{std::any} ma sens m.in. w luźno typowanych zasobnikach kontekstu,
gdy autor zasobnika nie zna listy możliwych typów, a odbiorca umie
jawnie odtworzyć typ właściwego wpisu. Jeśli rodzaje wartości są znane,
\texttt{variant} daje z reguły silniejszą kontrolę statyczną i czytelniejszą
obsługę przypadków. Unikaj \texttt{any} jako sposobu na ukrycie
źle określonych kontraktów interfejsu.

\section{\texttt{std::function}: wymazywanie typu obiektu wywoływalnego}
\label{sec:cpp14-function}

\texttt{std::function<R(Args...)>} (C++11, nagłówek
\texttt{<functional>}) przechowuje kopiowalny obiekt wywoływalny zgodny
z podaną sygnaturą. Może to być wskaźnik do funkcji, funktor,
wyrażenie lambda, a także obiekt związany z metodą składową. Konkretny
typ obiektu zostaje ukryty, ale interfejs jego wywołania jest znany:
\texttt{R(Args...)}.

\begin{lstlisting}[caption={Jeden typ std::function dla różnych obiektów wywoływalnych; kompletny program.}]
#include <cassert>
#include <functional>

int times_two(int x) { return 2 * x; }
struct PlusThree {
    int operator()(int x) const { return x + 3; }
};

int main() {
    std::function<int(int)> f = times_two;
    assert(f(4) == 8);

    f = PlusThree{};
    assert(f(4) == 7);

    int delta = 5;
    f = [delta](int x) { return x + delta; };
    assert(f(4) == 9);
}
\end{lstlisting}

\subsection{Pusty obiekt i koszt abstrakcji}

Domyślnie skonstruowany \texttt{std::function} jest pusty. Jego
\texttt{operator bool()} informuje, czy istnieje obiekt wywoływalny;
próba wywołania pustego obiektu zgłasza \texttt{std::bad\_function\_call}.
W porównaniu z bezpośrednim wywołaniem szablonu \texttt{std::function}
może wymagać pośredniego wywołania i dodatkowego przechowywania stanu,
a czasami także alokacji dynamicznej. Nie jest jednak poprawne twierdzenie,
że każdy \texttt{std::function} \emph{zawsze} wykonuje alokację.

\begin{lstlisting}[caption={Pusty std::function i jawne sprawdzanie; kompletny program.}]
#include <cassert>
#include <functional>

int main() {
    std::function<void()> callback;
    assert(!callback);
    callback = [] {};
    assert(static_cast<bool>(callback));
    callback();
    callback = nullptr;
    assert(!callback);
}
\end{lstlisting}

\subsection{Ograniczenie kopiowalności i czas życia przechwyceń}

W C++20 \texttt{std::function} wymaga kopiowalności typu przechowywanego
obiektu wywoływalnego. Dlatego lambda zawierająca \emph{move-only}
\texttt{std::unique\_ptr} nie może być bezpośrednio umieszczona w
\texttt{std::function}. W C++23 pojawił się
\texttt{std::move\_only\_function} dla takiego zastosowania.

Przechwytywanie przez referencję wymaga ponadto, aby obiekt odniesienia
pozostawał żywy podczas każdego późniejszego wywołania. Sam
\texttt{std::function} kontroluje czas życia \emph{przechowywanej lambdy},
ale nie przedłuża automatycznie czasu życia obiektów, do których lambda
odwołuje się przez referencje lub nieposiadające własności wskaźniki.

\begin{lstlisting}[caption={Callback utrzymujący obiekt przy życiu przez \texttt{shared\_ptr}; kompletny program.}]
#include <cassert>
#include <functional>
#include <memory>

int main() {
    auto state = std::make_shared<int>(10);
    std::function<int()> get_value = [state] {
        return *state;
    };
    state.reset();
    assert(get_value() == 10);
}
\end{lstlisting}

Trwałe przechwycenie \texttt{shared\_ptr} może z kolei utworzyć cykl
własności, np. gdy obiekt przechowuje callback, a callback przechwytuje
\texttt{shared\_ptr} do tego samego obiektu. Wówczas często lepiej
przechwycić \texttt{weak\_ptr} i uzyskać tymczasową własność przez
\texttt{lock()}.

\begin{lstlisting}[caption={Callback niewymuszający przedłużania życia obiektu; kompletny program.}]
#include <cassert>
#include <functional>
#include <memory>
#include <optional>

int main() {
    auto state = std::make_shared<int>(42);
    std::weak_ptr<int> weak = state;

    std::function<std::optional<int>()> read = [weak]()
        -> std::optional<int> {
        if (auto shared = weak.lock()) return *shared;
        return std::nullopt;
    };

    assert(read() == 42);
    state.reset();
    assert(!read().has_value());
}
\end{lstlisting}

\subsection{Szablon, wskaźnik do funkcji czy \texttt{std::function}?}

Jeżeli funkcja korzysta z obiektu wywoływalnego tylko podczas bieżącego
wywołania, często najprościej użyć szablonu. Pozwala to kompilatorowi
znać konkretny typ i daje okazję do \emph{inliningu}. Wskaźnik do funkcji
jest lekki, ale nie przechowa samodzielnie stanu przechwycenia lambdy.
\texttt{std::function} ułatwia natomiast przechowywanie i zastępowanie
heterogenicznych callbacków o wspólnej sygnaturze.

\begin{lstlisting}[caption={Wywołanie generyczne bez type erasure; kompletny program.}]
#include <cassert>
#include <functional>
#include <utility>

template<class F, class... Args>
decltype(auto) call_now(F&& f, Args&&... args) {
    return std::invoke(std::forward<F>(f),
                       std::forward<Args>(args)...);
}

int main() {
    int factor = 3;
    auto f = [factor](int x) { return factor * x; };
    assert(call_now(f, 7) == 21);
}
\end{lstlisting}

\section{Własny mechanizm type erasure: \texttt{Printable}}
\label{sec:cpp14-custom-erasure}

To rozszerzenie pokazuje, co może znajdować się pod gotowym
opakowaniem wymazującym typ. Przy pierwszej lekturze wystarczy
rozumieć, kiedy użyć \texttt{variant}, \texttt{any} lub
\texttt{function}; własną implementację można przeczytać później.

Pokażemy teraz, jak samodzielnie zbudować typ o stałym interfejsie,
który może przechowywać różne wartości drukowalne do strumienia. Celowo
nie wykorzystamy dziedziczenia typów użytkownika po wspólnej klasie
bazowej. Użytkownik powinien móc przekazać \texttt{int},
\texttt{std::string} albo własny typ, dla którego działa
\texttt{operator<<}. Wymagamy przy tym \emph{semantyki wartości}:
kopiowanie \texttt{Printable} ma kopiować przechowywany obiekt.

\paragraph{Idea implementacji.}
Wewnątrz opakowania znajdują się:
\begin{enumerate}
\item abstrakcyjny interfejs \texttt{Concept} z operacjami \texttt{print}
      i \texttt{clone};
\item szablonowy model \texttt{Model<T>}, który przechowuje konkretną
      wartość i realizuje interfejs;
\item wskaźnik \texttt{unique\_ptr<Concept>}, ukrywający konkretny
      typ w miejscu użycia \texttt{Printable}.
\end{enumerate}

\begin{lstlisting}[caption={Własne kopiowalne type erasure (C++20); kompletny program.}]
#include <cassert>
#include <concepts>
#include <memory>
#include <ostream>
#include <sstream>
#include <string>
#include <type_traits>
#include <utility>

class Printable {
    struct Concept {
        virtual ~Concept() = default;
        virtual void print(std::ostream&) const = 0;
        virtual std::unique_ptr<Concept> clone() const = 0;
    };

    template<class T>
    struct Model final : Concept {
        T value;
        template<class U>
        explicit Model(U&& x) : value(std::forward<U>(x)) {}
        void print(std::ostream& os) const override {
            os << value;
        }
        std::unique_ptr<Concept> clone() const override {
            return std::make_unique<Model<T>>(value);
        }
    };

    std::unique_ptr<Concept> self_;

public:
    Printable() = default;

    template<class T>
      requires (!std::same_as<std::remove_cvref_t<T>, Printable> &&
                std::copy_constructible<std::decay_t<T>> &&
                requires(std::ostream& os,
                         const std::decay_t<T>& value) {
                    { os << value } -> std::same_as<std::ostream&>;
                })
    Printable(T&& value)
        : self_(std::make_unique<Model<std::decay_t<T>>>(
              std::forward<T>(value))) {}

    Printable(const Printable& other)
        : self_(other.self_ ? other.self_->clone() : nullptr) {}

    Printable(Printable&&) noexcept = default;

    Printable& operator=(const Printable& other) {
        if (this != &other) {
            Printable tmp(other); // najpierw bezpieczna kopia
            swap(tmp);            // potem nie-rzucajaca zamiana
        }
        return *this;
    }

    Printable& operator=(Printable&&) noexcept = default;

    void swap(Printable& other) noexcept {
        self_.swap(other.self_);
    }

    void print(std::ostream& os) const {
        if (self_) self_->print(os);
        else os << "<empty>";
    }
};

int main() {
    Printable a = 17;
    Printable b = std::string("hello");
    Printable c = b; // gleboka kopia wartosci
    b = 7.5;

    std::ostringstream stream;
    a.print(stream);
    c.print(stream);
    assert(stream.str() == "17hello");
}
\end{lstlisting}

\subsection{Analiza własności i kosztów}

Ten interfejs jest \emph{otwarty}: każdy nowy typ spełniający wymagania
może zostać opakowany bez modyfikowania klasy \texttt{Printable}.
Użytkownik nie musi dziedziczyć po klasie bazowej, co jest ważną zaletą
w porównaniu ze zwykłym polimorfizmem opartym na \texttt{Base}.
\texttt{Model<T>} jest instancjonowany w czasie kompilacji, natomiast
\texttt{print()} jest wywoływane dynamicznie przez ukryty interfejs.
Kopiowanie przez \texttt{clone()} zachowuje rzeczywisty typ wartości.

Wersja edukacyjna wykorzystuje osobną alokację dynamiczną przy
przechowywaniu niepustej wartości i kolejną alokację przy głębokim
kopiowaniu. Nie wykorzystuje \emph{small-buffer optimization}, które
mogłoby przechowywać niewielkie obiekty bezpośrednio wewnątrz opakowania.
W produkcyjnym type erasure trzeba dodatkowo przemyśleć \emph{wyrównanie},
przenoszenie, bezpieczne niszczenie, usuwanie typów niekopiowalnych,
kwalifikatory \texttt{const}, obsługę wyjątków i stabilność ABI.

\paragraph{Gwarancja silna przypisania kopiującego.}
Tworzenie obiektu \texttt{tmp(other)} może zakończyć się wyjątkiem.
Wówczas obiekt docelowy pozostaje niezmieniony. Jeśli kopia się uda,
\texttt{swap(tmp)} jest operacją \texttt{noexcept}; stara zawartość
zostanie zniszczona razem z \texttt{tmp}. Jest to klasyczny idiom
\emph{copy-and-swap}. Alternatywą jest niekopiowalne opakowanie,
które usuwa operacje kopiowania i przechowuje tylko obiekty przenoszalne.

\section{Type erasure a dziedziczenie i polimorfizm statyczny}
\label{sec:cpp14-comparison}

\subsection{Wspólna baza klas}

Gdy typy zostały zaprojektowane jako część wspólnej hierarchii,
\texttt{std::unique\_ptr<Base>} pozwala przechowywać obiekty o różnych
typach dynamicznych i wywoływać ich operacje wirtualne. Implementacje
muszą dziedziczyć po \texttt{Base}, a tworzenie kopii polimorficznej
wymaga własnego protokołu (np. \texttt{clone()}). Z drugiej strony
\texttt{unique\_ptr} jasno wyraża własność i nie kopiuje obiektów
wskutek zwykłego kopiowania uchwytu (ponieważ kopiowanie jest zabronione).

\begin{lstlisting}[caption={Polimorfizm przez bazę i \texttt{unique\_ptr}; kompletny program.}]
#include <cassert>
#include <memory>
#include <string>
#include <vector>

struct Shape {
    virtual ~Shape() = default;
    virtual int area() const = 0;
};
struct Square final : Shape {
    int side;
    explicit Square(int s) : side(s) {}
    int area() const override { return side * side; }
};
struct Rectangle final : Shape {
    int width, height;
    Rectangle(int w, int h) : width(w), height(h) {}
    int area() const override { return width * height; }
};

int main() {
    std::vector<std::unique_ptr<Shape>> shapes;
    shapes.push_back(std::make_unique<Square>(3));
    shapes.push_back(std::make_unique<Rectangle>(2, 5));
    assert(shapes[0]->area() + shapes[1]->area() == 19);
}
\end{lstlisting}

\subsection{Kiedy nie wymazywać typu?}

Jeżeli algorytm potrzebuje konkretnego typu wyłącznie podczas jednego
wywołania, parametr szablonowy z \texttt{concept} jest często prostszy
i pozwala unikać alokacji. Type erasure ma sens przede wszystkim na
granicy interfejsów, gdy wartości o różnych typach muszą być
\emph{przechowywane w jednej strukturze}, przekazywane przez stabilny
interfejs albo zastępowane w trakcie działania programu.

\begin{center}
\small
\begin{tabularx}{\textwidth}{@{}lXX@{}}
\toprule
\textbf{Rozwiązanie} & \textbf{Zalety} & \textbf{Ograniczenia} \\
\midrule
\texttt{variant<Ts...>} & Typy znane statycznie, sprawdzanie przypadków, lokalne przechowywanie wartości & Zamknięty katalog typów, złożoność odwiedzania i rozmiar największej alternatywy \\
\texttt{any} & Dowolna kopiowalna wartość, wspólny typ pojemnika & Trzeba odzyskać dokładny typ; brak wspólnego interfejsu operacji \\
\texttt{function} & Jedna sygnatura dla różnych callbacków & Kopiowalność (C++20), pośrednictwo, możliwe alokacje \\
Własne type erasure & Otwarty katalog modeli bez wspólnej bazy typów klienta & Złożona implementacja, zarządzanie pamięcią i wyjątkami \\
\texttt{unique\_ptr<Base>} & Jawna hierarchia, dziedziczone operacje, kontrola własności & Wymaga wspólnej bazy; operacje wirtualne \\
Szablon z \texttt{concept} & Typ znany kompilatorowi, możliwość optymalizacji & Wiele instancjacji; zwykle nie przechowa różnych typów w jednej kolekcji \\
\bottomrule
\end{tabularx}
\end{center}

\section{Wydajność: układ pamięci i koszty dynamicznej dyspozycji}
\label{sec:cpp14-performance}

Nie ma uniwersalnej odpowiedzi na pytanie, który mechanizm jest
najszybszy. Poniżej wymieniamy dominujące koszty; faktyczne zachowanie
zależy od rozmiaru typów, optymalizacji kompilatora, struktury programu
i rozkładu typów w danych.

\begin{itemize}
\item \texttt{optional<T>} dodaje znacznik obecności; dostęp do zawartości
      może wymagać sprawdzenia warunku.
\item \texttt{variant<Ts...>} przechowuje wystarczającą ilość miejsca
      na największy typ i indeks aktywnej alternatywy; odwiedzający
      może być implementowany przez instrukcję wyboru lub tablicę
      funkcji, a szczegóły zależą od kompilatora.
\item \texttt{any} przechowuje metadane typu oraz wartość, potencjalnie
      w osobno alokowanym buforze.
\item \texttt{function} może wymagać pośredniego wywołania i przechowania
      stanu lambdy; czasami da się uniknąć alokacji.
\item Własna implementacja z \texttt{unique\_ptr<Concept>} wykonuje
      alokację dynamiczną i pośrednie wywołanie wirtualne.
\end{itemize}

Przede wszystkim trzeba mierzyć \emph{rzeczywisty} przypadek użycia.
Jednorazowy koszt konstruowania callbacku może być nieistotny, gdy jego
wykonanie trwa długo, ale może dominować w pętli składającej się z
milionów krótkich operacji.

\subsection{Test bez mylenia czasu kompilacji z wykonaniem}

Rzetelne porównanie wymaga rozgrzania kodu, wykorzystania odpowiedniego
zestawu danych, uwzględnienia optymalizacji i sprawdzenia, czy kompilator
nie usunął mierzonej pracy. Warto rozdzielać koszt \emph{tworzenia}
opakowania od kosztu \emph{wywołania} i analizować nie tylko średni
czas, ale też rozkład wyników oraz liczbę alokacji. Benchmarku
porównującego dwa mechanizmy nie traktujemy jako dowodu ich przewagi
we wszystkich zastosowaniach.

\section{Granice bezpieczeństwa: wyjątki, czas życia i współbieżność}

\paragraph{Wyjątki.}
\texttt{optional::value()}, \texttt{std::get} dla wariantu, \texttt{any\_cast}
i wywołanie pustego \texttt{std::function} mogą zgłaszać odpowiednie
wyjątki. Konstrukcja zawartości i przypisanie także mogą rzucać,
ponieważ wykonują operacje typu przechowywanego. Należy rozróżniać
\emph{brak wartości} (normalny stan \texttt{optional}) od
\emph{niepowodzenia operacji}, któremu warto nadać osobną reprezentację.

\paragraph{Czas życia.}
\texttt{variant}, \texttt{optional} i \texttt{any} posiadają własne
przechowywane wartości. Można jednak umieścić w nich wskaźniki,
\texttt{reference\_wrapper} albo inne nieposiadające własności uchwyty;
wtedy odpowiedzialność za czas życia wskazywanych danych pozostaje po
stronie programisty. To samo dotyczy referencji przechwyconych przez
lambdę umieszczoną w \texttt{std::function}.

\paragraph{Wątki.}
Żaden z opisanych typów nie czyni dowolnych współbieżnych modyfikacji
tego samego obiektu bezpiecznymi. Przykładowo równoczesny zapis do
jednego \texttt{variant} i odczyt jego aktywnej zawartości bez
synchronizacji mogą powodować wyścig danych. Gdy callback jest
uruchamiany równolegle, trzeba osobno określić zasady synchronizacji
stanu przechwyconego przez lambdę.

\section{Nowości po C++20 i świadomy wybór standardu}

W całym rozdziale przykłady są zgodne z C++20. W nowszych standardach
pojawiają się udogodnienia, które warto znać, ale których \emph{nie wolno}
zakładać przy kompilacji z \texttt{-std=c++20}:
\begin{itemize}
\item \textbf{C++23:} \texttt{std::expected<T,E>} --- wynik albo
      jawna wartość błędu;
\item \textbf{C++23:} \texttt{std::move\_only\_function} --- opakowanie
      dla obiektów wywoływalnych, które nie muszą być kopiowalne;
\item \textbf{C++23:} operacje monadyczne dla \texttt{std::optional}:
      \texttt{and\_then}, \texttt{transform} i \texttt{or\_else};
\item \textbf{C++26:} \texttt{std::function\_ref} --- niewłaścicielski
      widok na obiekt wywoływalny, którego czas życia musi zostać
      zapewniony przez wywołującego.
\end{itemize}

Dostępność funkcji bibliotecznych w konkretnym kompilatorze i jego
bibliotece standardowej może być opóźniona w stosunku do daty
opublikowania standardu.

\section{Zasady projektowania i typowe błędy}
\label{sec:cpp14-guidelines}

\begin{enumerate}
\item Wybieraj \texttt{optional<T>}, gdy brak wyniku jest normalny i
      nie potrzebujesz szczegółowego typu błędu.
\item Wybieraj \texttt{variant}, gdy zestaw przypadków jest znany
      i zamknięty; obsługuj je przez \texttt{visit}, jeśli pozwala to
      uniknąć łańcucha ręcznych kontroli typu.
\item Stosuj \texttt{any} tylko tam, gdzie świadomie rezygnujesz
      ze statycznie określonego zestawu typów i masz jasny protokół
      odzyskiwania ich wartości.
\item Używaj \texttt{std::function} do \emph{przechowywanych},
      kopiowalnych callbacków; dla funkcji wykonywanej natychmiast
      rozważ parametr szablonowy.
\item Nie zakładaj, że \texttt{std::function} zawsze alokuje ani
      że \texttt{std::variant} zawsze jest szybszy. Zmierz przypadek
      zastosowania.
\item Nie pozostawiaj referencji lub wskaźników do obiektów, których
      czas życia zakończy się przed użyciem przechowywanego callbacku.
\item Unikaj \texttt{std::get} i \texttt{any\_cast} w sytuacjach,
      w których nie masz gwarancji typu; do sprawdzania wykorzystuj
      \texttt{get\_if} i wskaźnikowy \texttt{any\_cast}.
\item Nie używaj \texttt{std::any} do bezpośredniego opakowania
      typów przenoszalnych, ale niekopiowalnych, np. \texttt{unique\_ptr}.
\item Rozdziel decyzję o \emph{otwartym zbiorze implementacji} od
      decyzji o \emph{dynamicznym czasie życia i alokacji}.
\end{enumerate}

\section*{Podsumowanie rozdziału}
\addcontentsline{toc}{section}{Podsumowanie rozdziału 10}

\texttt{optional} modeluje brak lub obecność wartości, a
\texttt{variant} modeluje zamkniętą sumę przypadków. \texttt{any}
ukrywa typ kopiowalnej wartości, nie narzucając jej operacji.
\texttt{function} ukrywa typ obiektu wywoływalnego przy zachowaniu
ustalonej sygnatury. Własne type erasure pozwala opakować otwartą
rodzinę obcych typów bez wymagania wspólnej klasy bazowej, ale wymaga
starannego określenia semantyki własności, kopiowania i wyjątków.
Dobry projekt zaczyna się więc od określenia kontraktu: \emph{które
przypadki istnieją, kto posiada obiekty, jakie operacje udostępniamy
oraz czy zbiór typów ma pozostać otwarty}.

\section*{Dokumentacja i literatura}
\addcontentsline{toc}{section}{Dokumentacja i literatura rozdziału 10}
\begin{itemize}
\item \url{https://en.cppreference.com/w/cpp/utility/optional} --- \texttt{std::optional}.
\item \url{https://en.cppreference.com/w/cpp/utility/variant} --- \texttt{std::variant}.
\item \url{https://en.cppreference.com/w/cpp/utility/any} --- \texttt{std::any}.
\item \href{https://en.cppreference.com/w/cpp/utility/functional/function}{cppreference: \texttt{std::function}} --- dokumentacja opakowania wywołań.
\item Bjarne Stroustrup, \emph{A Tour of C++}, wyd. 3, Addison--Wesley, 2022.
\item Klaus Iglberger, \emph{C++ Software Design}, O'Reilly, 2022.
\end{itemize}

% Kompletne programy sa oznaczone w podpisach listingow.

\chapter{Wymazywanie typów i typy sumaryczne}
\label{chap:cpp14-type-erasure}

\section{Problem: jak reprezentować różne rodzaje wartości?}

W programowaniu generycznym często chcemy posługiwać się różnymi typami
przez wspólny interfejs. C++ oferuje kilka mechanizmów realizacji tego celu,
ale mechanizmy te nie są zamienne.
Wyobraźmy sobie aplikację z trzema poleceniami: zapisaniem pliku,
dodaniem dwóch liczb i zakończeniem pracy. Gdy znamy wszystkie trzy
rodzaje poleceń, możemy je wymienić w \texttt{std::variant}. Jeśli
zewnętrzne moduły mają dodawać nowe polecenia bez zmiany tej listy,
potrzebujemy otwartego interfejsu, np. klasy bazowej z funkcją wirtualną.
Wybór mechanizmu zaczynamy więc od pytań:
\begin{enumerate}
\item Czy zestaw dozwolonych typów jest \emph{zamknięty} (ustalony z góry),
      czy \emph{otwarty} (może zostać rozszerzony o nowy typ)?
\item Czy chcemy przechować \emph{wartość}, czy jedynie referencję/uchwyt
      do istniejącego obiektu?
\item Czy wywoływane operacje mają być znane już podczas kompilacji,
      czy wybierane dynamicznie w czasie działania programu?
\item Czy dopuszczamy alokację dynamiczną, kopiowanie oraz koszt
      pośredniego wywołania funkcji?
\end{enumerate}

\emph{Typ sumaryczny} (ang. \emph{sum type}) reprezentuje alternatywę:
wartość należy do jednego z kilku przypadków. W C++17 taką rolę pełni
\texttt{std::variant<Ts...>}; \texttt{std::optional<T>} opisuje wariant
,,brak wartości albo wartość \texttt{T}''. Natomiast \emph{type erasure}
(\emph{wymazywanie typu}) pozwala ukryć konkretny typ obiektu za
interfejsem nieparametryzowanym przez ten typ. Standardowe przykłady
to \texttt{std::any} i \texttt{std::function}.

\begin{center}
\small
\begin{tabularx}{\textwidth}{@{}lXX@{}}
\toprule
\textbf{Mechanizm} & \textbf{Informacja o typie} & \textbf{Zasadnicze zastosowanie} \\
\midrule
\texttt{optional<T>} & Jedna możliwość \texttt{T} lub brak & Wynik opcjonalny, wartość nieobecna \\
\texttt{variant<Ts...>} & Zamknięty zbiór alternatyw & Zdarzenia, komendy, model stanu \\
\texttt{any} & Typ ustalany w czasie wykonania & Heterogeniczne wartości bez wspólnego interfejsu operacji \\
\texttt{function<R(Args...)>} & Dowolny zgodny obiekt wywoływalny & Przechowywany callback, funkcja jako wartość \\
\texttt{unique\_ptr<Base>} & Otwarta hierarchia klas & Polimorfizm dynamiczny i tożsamość obiektu \\
Szablon/\texttt{concept} & Typ znany przy instancjowaniu & Polimorfizm statyczny, specjalizacja kodu \\
\bottomrule
\end{tabularx}
\end{center}

Najprostsza ściąga to: \texttt{optional<int>} oznacza liczbę albo jej
brak, \texttt{variant<int, std::string>} jedną z dwóch znanych
postaci wartości, a \texttt{any} wartość o typie, którego autor
pojemnika nie musi znać. \texttt{function<void()>} przechowuje
różne działania, które można wywołać w taki sam sposób. W kolejnych
sekcjach zobaczymy, co z tych wyborów wynika dla sprawdzania typów,
własności danych i kosztu wykonania.

\section{\texttt{std::optional<T>}: wartość, której może nie być}

\begin{center}\small
\begin{tabularx}{\textwidth}{@{}p{0.33\textwidth}X@{}}
\toprule
Typ lub operacja & Co reprezentuje \\
\midrule
\code{std::optional<T>} & Wartość typu \code{T} albo jej brak; nie opisuje przyczyny błędu. \\
\code{has_value()} & Sprawdza, czy \code{optional} przechowuje wartość. \\
\code{value()} & Zwraca wartość lub zgłasza \code{std::bad_optional_access}, gdy jej brak. \\
\code{value_or(fallback)} & Zwraca wartość albo podany zamiennik, konwertowany do \code{T}. \\
\code{std::variant<Ts...>} & Dokładnie jedną z ustalonych alternatyw typów. \\
\code{std::any} & Wartość typu ustalonego dopiero w czasie wykonania; przy odczycie trzeba sprawdzić typ. \\
\bottomrule
\end{tabularx}
\end{center}
\label{sec:cpp14-optional}

Typ \texttt{std::optional<T>} (C++17, nagłówek \texttt{<optional>})
przechowuje bezpośrednio obiekt typu \texttt{T}, jeżeli jest on obecny,
lub reprezentuje stan pusty. Jest to pełnoprawna \emph{wartość}, a nie
wskaźnik. W szczególności pusty \texttt{optional} nie oznacza wskaźnika
zerowego, a niepusty nie wymaga automatycznie osobnej alokacji dynamicznej.

\begin{lstlisting}[caption={Najważniejsze operacje optional (C++17); kompletny program.}]
#include <cassert>
#include <optional>
#include <string>

int main() {
    std::optional<int> x;         // stan pusty
    assert(!x.has_value());
    assert(x == std::nullopt);

    x = 42;
    assert(x.has_value());
    assert(*x == 42);
    assert(x.value() == 42);
    assert(x.value_or(-1) == 42);

    x.reset();
    assert(x.value_or(-1) == -1);
    x.emplace(7);
    assert(*x == 7);

    std::optional<std::string> name{std::in_place, 3, 'a'};
    assert(*name == "aaa");
}
\end{lstlisting}

Dostęp przez \texttt{*x} i \texttt{x->...} wymaga obecności wartości:
\emph{użycie pustego obiektu w taki sposób jest niezdefiniowanym
zachowaniem}. Metoda \texttt{value()} wykonuje sprawdzenie i dla stanu
pustego zgłasza \texttt{std::bad\_optional\_access}. Warto pamiętać,
że \texttt{value\_or(default)} przyjmuje wartość domyślną jako argument
funkcji: wyrażenie tworzące argument jest obliczane również wtedy,
gdy \texttt{optional} zawiera wartość. Może ponadto następować kopiowanie
lub przeniesienie przechowywanego obiektu.

\subsection{Praktyczny przykład: parsowanie liczby}

Funkcja parsująca dane wejściowe może zwrócić \texttt{optional}, jeśli
interesuje nas tylko informacja ,,udało się/nie udało się''.
\texttt{std::from\_chars} unika wyjątków i nie wymaga tworzenia dodatkowego
łańcucha zakończonego zerem. W przykładzie wymagamy, aby \emph{cały}
łańcuch był prawidłową liczbą.

\begin{lstlisting}[caption={Parsowanie int do optional (C++17); kompletny program.}]
#include <cassert>
#include <charconv>
#include <optional>
#include <string_view>
#include <system_error>

std::optional<int> parse_int(std::string_view text) {
    int result = 0;
    auto [ptr, ec] = std::from_chars(
        text.data(), text.data() + text.size(), result);
    if (ec != std::errc{} || ptr != text.data() + text.size())
        return std::nullopt;
    return result;
}

int main() {
    assert(parse_int("123") == 123);
    assert(parse_int("123x") == std::nullopt);
    assert(parse_int("") == std::nullopt);
    assert(parse_int("99999999999999999999") == std::nullopt);
}
\end{lstlisting}

\subsection{Własność, referencje i koszt}

\texttt{optional<T>} posiada swoją wartość i zwykle zajmuje więcej pamięci
niż sam \texttt{T}, ponieważ musi przechowywać informację, czy wartość
istnieje. Rozmiar zależy od typu i implementacji. Standard C++20
nie udostępnia \texttt{std::optional<T\&>}. Jeśli wynik ma wskazywać na
istniejący obiekt, można użyć wskaźnika nieposiadającego własności albo
\texttt{std::optional<std::reference\_wrapper<T>>}. W obu przypadkach
programista odpowiada za czas życia wskazywanego obiektu.

\begin{lstlisting}[caption={Opcjonalna referencja: brak przenoszenia własności; fragment.}]
#include <functional>
#include <optional>
#include <string>

std::optional<std::reference_wrapper<const std::string>>
find_name(bool exists, const std::string& name) {
    if (exists) return std::cref(name);
    return std::nullopt;
}
// Zwracana referencja jest poprawna tylko gdy `name` nadal istnieje.
\end{lstlisting}

Gdy brak wartości jest sytuacją normalną, \texttt{optional} najczęściej
czytelniej wyraża intencję niż wyjątkowe zakończenie funkcji. Jeżeli
konieczna jest informacja \emph{dlaczego} nie uzyskano wyniku, należy
zaprojektować osobny typ błędu, użyć \texttt{variant} lub (w C++23)
\texttt{std::expected<T,E>}.

\section{\texttt{std::variant<Ts...>}: zamknięta alternatywa typów}
\label{sec:cpp14-variant}

\texttt{std::variant<Ts...>} (C++17, nagłówek \texttt{<variant>})
przechowuje \emph{jedną aktywną alternatywę} z listy typów \texttt{Ts...}.
Nie jest odpowiednikiem \texttt{std::tuple<Ts...>}: krotka zawiera
\emph{wszystkie} składowe, a wariant dokładnie jeden z dopuszczonych
przypadków (poza szczególnym stanem \texttt{valueless\_by\_exception()}).

\begin{lstlisting}[caption={Tworzenie i badanie aktywnej alternatywy (C++17); kompletny program.}]
#include <cassert>
#include <string>
#include <variant>

int main() {
    using Value = std::variant<int, std::string>;
    Value x = 7;
    assert(x.index() == 0);
    assert(std::holds_alternative<int>(x));
    assert(std::get<int>(x) == 7);

    x = std::string("hello");
    assert(x.index() == 1);
    assert(std::get_if<int>(&x) == nullptr);
    assert(std::get_if<std::string>(&x) != nullptr);

    x.emplace<0>(19);
    assert(std::get<0>(x) == 19);
}
\end{lstlisting}

Niepoprawne \texttt{std::get<T>(x)} (gdy aktywny jest inny typ)
zgłasza \texttt{std::bad\_variant\_access}. Alternatywa
\texttt{std::get\_if<T>(\&x)} zwraca \texttt{nullptr}, a zatem nadaje się
do sprawdzania bez wyjątku. \texttt{std::get<T>} i
\texttt{std::holds\_alternative<T>} wymagają, aby \texttt{T} występował
\emph{dokładnie raz} w liście alternatyw; dla powtarzających się typów
używamy indeksów. Nie należy także zapominać o możliwym niepowodzeniu
operacji konwersji przy konstruowaniu wariantu.

\subsection{Domyślna alternatywa i \texttt{std::monostate}}

Domyślnie konstruowany wariant próbuje skonstruować pierwszą alternatywę.
Jeżeli chcemy jawnie reprezentować stan pusty, można umieścić na początku
\texttt{std::monostate}; zapewnia on także domyślną konstruowalność
wariantu, gdy inny pierwszy typ nie jest domyślnie konstruowalny.

\begin{lstlisting}[caption={Jawny stan początkowy monostate; kompletny program.}]
#include <cassert>
#include <string>
#include <variant>

int main() {
    using State = std::variant<std::monostate, int, std::string>;
    State s;
    assert(std::holds_alternative<std::monostate>(s));
    s = std::string("ready");
    assert(std::get<std::string>(s) == "ready");
}
\end{lstlisting}

\subsection{Odwiedzanie wariantu: \texttt{std::visit}}

Samo wielokrotne sprawdzanie \texttt{holds\_alternative} szybko
prowadzi do rozbudowanych instrukcji \texttt{if}. Mechanizm
\texttt{std::visit(visitor, variant)} wywołuje odwiedzającego z aktywną
wartością odpowiedniego typu. Gdy odwiedzający jest lambdą generyczną,
kompilator instancjuje jej operator wywołania dla różnych alternatyw.

\begin{lstlisting}[caption={Jeden odwiedzający generyczny; kompletny program.}]
#include <cassert>
#include <cstddef>
#include <string>
#include <type_traits>
#include <variant>

int main() {
    std::variant<int, double, std::string> v = "hi";

    const std::size_t size = std::visit([](const auto& x) {
        using T = std::decay_t<decltype(x)>;
        if constexpr (std::is_same_v<T, std::string>)
            return x.size();
        else
            return sizeof(x);
    }, v);

    assert(size == 2);
}
\end{lstlisting}

Do osobnej obsługi poszczególnych typów wygodny jest wzorzec
\texttt{overloaded}, znany również z rozdziału o lambdach.
Wykorzystuje on dziedziczenie po kilku obiektach funkcyjnych oraz
ujawnienie wszystkich operatorów wywołania za pomocą \texttt{using}.

\begin{lstlisting}[caption={overloaded oraz std::visit (C++17); kompletny program.}]
#include <cassert>
#include <cstddef>
#include <string>
#include <variant>

template<class... Fs>
struct overloaded : Fs... {
    using Fs::operator()...;
};
template<class... Fs> overloaded(Fs...) -> overloaded<Fs...>;

int main() {
    using Data = std::variant<int, double, std::string>;
    Data data = std::string("abc");

    auto length = std::visit(overloaded{
        [](int) { return std::size_t{1}; },
        [](double) { return std::size_t{1}; },
        [](const std::string& s) { return s.size(); }
    }, data);
    assert(length == 3);
}
\end{lstlisting}

W przypadku odwiedzania zwracającego wartość wymagany jest zgodny typ
wyniku wszystkich możliwych wywołań (w praktyce należy zapewnić
identyczny typ i kategorię wartości zwracanych wyrażeń, z wyjątkiem
szczególnych reguł dla \texttt{void}). Przy wielu wariantach
\texttt{std::visit} uwzględnia \emph{kombinacje} alternatyw. To daje
czytelny kod przy małej liczbie przypadków, ale może zwiększyć liczbę
instancjacji szablonów.

\begin{lstlisting}[caption={Odwiedzenie dwóch wariantów naraz; kompletny program.}]
#include <cassert>
#include <variant>

int main() {
    std::variant<int, double> a = 3;
    std::variant<int, double> b = 2.5;
    const double sum = std::visit(
        [](const auto& x, const auto& y) -> double {
            return static_cast<double>(x) + static_cast<double>(y);
        }, a, b);
    assert(sum == 5.5);
}
\end{lstlisting}

\subsection{Modyfikacja stanu i kwalifikacja \texttt{const}}

\texttt{std::visit} respektuje typ i kategorię wartości przekazanego
wariantu. Odwiedzając obiekt przez \texttt{variant\&}, możemy modyfikować
aktywną alternatywę; przez \texttt{const variant\&} uzyskujemy
referencję stałą. Przy przekazaniu \texttt{std::move(v)} odwiedzający
może otrzymać referencję rvalue i przenieść przechowywaną wartość.
Przeniesienie nie musi zmieniać aktywnego indeksu wariantu, ale zawartość
przeniesionego obiektu może zostać zmieniona zgodnie z semantyką typu.

\begin{lstlisting}[caption={Modyfikacja aktywnej alternatywy; kompletny program.}]
#include <cassert>
#include <variant>

int main() {
    std::variant<int, double> v = 10;
    std::visit([](auto& x) { x += 1; }, v);
    assert(std::get<int>(v) == 11);
}
\end{lstlisting}

\subsection{Stan \texttt{valueless\_by\_exception}}

Wyjątek podczas przypisania lub tworzenia nowej wartości może
pozostawić \texttt{variant} bez aktywnej alternatywy. Ten stan
wykrywa metoda \texttt{valueless\_by\_exception()}.
Wtedy \texttt{index()} zwraca \texttt{std::variant\_npos}.
Wywołanie \texttt{std::visit} zgłasza \texttt{std::bad\_variant\_access}.
Nie jest to zwykła ,,pusta'' alternatywa. Jeśli jest potrzebna,
można użyć \texttt{std::monostate}.

\paragraph{Ważne ograniczenia.}
\texttt{variant} nie przechowuje typów referencyjnych ani \texttt{void};
nie można także dodać dowolnego nowego typu bez zmiany listy alternatyw.
Duży wariant może zajmować pamięć odpowiadającą największej alternatywie
plus informację o aktywnym przypadku i ewentualne wyrównanie; dokładny
rozmiar nie jest stałą wielkością określoną przez standard.

\section{Przykład projektowy: interpreter prostych poleceń}
\label{sec:cpp14-commands}

Rozpatrzmy miniaturowy model poleceń. Każdy rodzaj polecenia ma inne pola,
ale wszystkie należą do zamkniętej rodziny. \texttt{variant} oddziela
\emph{dane} od \emph{operacji}, a \texttt{visit} pozwala jednoznacznie
obsłużyć każdy przypadek.

\begin{lstlisting}[caption={Modelowanie poleceń za pomocą variant; kompletny program.}]
#include <cassert>
#include <string>
#include <variant>
#include <vector>

struct Add  { int a, b; };
struct Save { std::string filename; };
struct Quit {};
using Command = std::variant<Add, Save, Quit>;

struct Execute {
    int& total;
    bool& stopped;

    void operator()(const Add& c) const {
        total += c.a + c.b;
    }
    void operator()(const Save& c) const {
        if (c.filename.empty()) stopped = true;
    }
    void operator()(const Quit&) const {
        stopped = true;
    }
};

int main() {
    std::vector<Command> commands{
        Add{2, 3}, Save{"result.txt"}, Add{4, 5}, Quit{}
    };
    int total = 0;
    bool stopped = false;
    for (const auto& command : commands) {
        if (stopped) break;
        std::visit(Execute{total, stopped}, command);
    }
    assert(total == 14);
    assert(stopped);
}
\end{lstlisting}

Dodanie nowego typu polecenia, np. \texttt{Delete}, wymaga rozszerzenia
\texttt{Command} i obsługi nowego przypadku. W dobrze skonstruowanym
odwiedzającym brakująca obsługa będzie błędem kompilacji. To zaleta
typów sumarycznych przy modelowaniu \emph{zamkniętego} katalogu zdarzeń.

\section{\texttt{std::any}: wartość o nieznanym typie}
\label{sec:cpp14-any}

\texttt{std::any} (C++17, nagłówek \texttt{<any>}) może przechowywać
obiekt dowolnego typu spełniającego wymaganie kopiowalności i pozwala
później odzyskać dokładny typ przez \texttt{std::any\_cast}. To klasyczna
implementacja \emph{type erasure}: użytkownik obiektu \texttt{any}
nie zna szablonowego typu przechowywanej wartości.

\begin{lstlisting}[caption={std::any i bezpieczne wydobycie typu (C++17); kompletny program.}]
#include <any>
#include <cassert>
#include <string>

int main() {
    std::any value = 17;
    assert(value.has_value());
    assert(std::any_cast<int>(value) == 17);

    if (auto ptr = std::any_cast<int>(&value)) {
        *ptr += 1;
    }
    assert(std::any_cast<int>(value) == 18);

    value.emplace<std::string>("text");
    assert(std::any_cast<int>(&value) == nullptr);
    assert(std::any_cast<const std::string&>(value) == "text");

    value.reset();
    assert(!value.has_value());
}
\end{lstlisting}

\texttt{std::any\_cast<T>(value)} zgłasza \texttt{std::bad\_any\_cast},
jeśli aktywny typ nie odpowiada żądanemu. Wariant wskaźnikowy
\texttt{std::any\_cast<T>(\&value)} zwraca \texttt{nullptr}; jest zatem
wygodny do ostrożnego sprawdzania typu. \texttt{type()} udostępnia
\texttt{std::type\_info} opisujący aktualnie przechowywany typ,
a dla pustego obiektu zwraca \texttt{typeid(void)}.

\begin{lstlisting}[caption={\texttt{any\_cast} nie jest konwersją numeryczną; kompletny program.}]
#include <any>
#include <cassert>
#include <typeinfo>

int main() {
    std::any value = 42;
    assert(value.type() == typeid(int));
    assert(std::any_cast<double>(&value) == nullptr);
    assert(std::any_cast<int>(&value) != nullptr);
}
\end{lstlisting}

\texttt{any} przechowuje \emph{wartość}, nie udostępnia jednak uniwersalnego
operatora \texttt{+}, porównania czy drukowania tej wartości. Aby wykonać
operację, trzeba znać typ albo dostarczyć odrębny mechanizm dyspozycji.
Ponadto przechowywany typ musi być kopiowalny: \texttt{std::any}
nie nadaje się do bezpośredniego przechowywania \texttt{std::unique\_ptr<T>}.
Niektóre implementacje wykorzystują \emph{small-object optimization},
ale standard nie obiecuje uniknięcia alokacji dla konkretnego małego typu.

\paragraph{Kiedy używać?}
\texttt{std::any} ma sens m.in. w luźno typowanych zasobnikach kontekstu,
gdy autor zasobnika nie zna listy możliwych typów, a odbiorca umie
jawnie odtworzyć typ właściwego wpisu. Jeśli rodzaje wartości są znane,
\texttt{variant} daje z reguły silniejszą kontrolę statyczną i czytelniejszą
obsługę przypadków. Unikaj \texttt{any} jako sposobu na ukrycie
źle określonych kontraktów interfejsu.

\section{\texttt{std::function}: wymazywanie typu obiektu wywoływalnego}
\label{sec:cpp14-function}

\texttt{std::function<R(Args...)>} (C++11, nagłówek
\texttt{<functional>}) przechowuje kopiowalny obiekt wywoływalny zgodny
z podaną sygnaturą. Może to być wskaźnik do funkcji, funktor,
wyrażenie lambda, a także obiekt związany z metodą składową. Konkretny
typ obiektu zostaje ukryty, ale interfejs jego wywołania jest znany:
\texttt{R(Args...)}.

\begin{lstlisting}[caption={Jeden typ std::function dla różnych obiektów wywoływalnych; kompletny program.}]
#include <cassert>
#include <functional>

int times_two(int x) { return 2 * x; }
struct PlusThree {
    int operator()(int x) const { return x + 3; }
};

int main() {
    std::function<int(int)> f = times_two;
    assert(f(4) == 8);

    f = PlusThree{};
    assert(f(4) == 7);

    int delta = 5;
    f = [delta](int x) { return x + delta; };
    assert(f(4) == 9);
}
\end{lstlisting}

\subsection{Pusty obiekt i koszt abstrakcji}

Domyślnie skonstruowany \texttt{std::function} jest pusty. Jego
\texttt{operator bool()} informuje, czy istnieje obiekt wywoływalny;
próba wywołania pustego obiektu zgłasza \texttt{std::bad\_function\_call}.
W porównaniu z bezpośrednim wywołaniem szablonu \texttt{std::function}
może wymagać pośredniego wywołania i dodatkowego przechowywania stanu,
a czasami także alokacji dynamicznej. Nie jest jednak poprawne twierdzenie,
że każdy \texttt{std::function} \emph{zawsze} wykonuje alokację.

\begin{lstlisting}[caption={Pusty std::function i jawne sprawdzanie; kompletny program.}]
#include <cassert>
#include <functional>

int main() {
    std::function<void()> callback;
    assert(!callback);
    callback = [] {};
    assert(static_cast<bool>(callback));
    callback();
    callback = nullptr;
    assert(!callback);
}
\end{lstlisting}

\subsection{Ograniczenie kopiowalności i czas życia przechwyceń}

W C++20 \texttt{std::function} wymaga kopiowalności typu przechowywanego
obiektu wywoływalnego. Dlatego lambda zawierająca \emph{move-only}
\texttt{std::unique\_ptr} nie może być bezpośrednio umieszczona w
\texttt{std::function}. W C++23 pojawił się
\texttt{std::move\_only\_function} dla takiego zastosowania.

Przechwytywanie przez referencję wymaga ponadto, aby obiekt odniesienia
pozostawał żywy podczas każdego późniejszego wywołania. Sam
\texttt{std::function} kontroluje czas życia \emph{przechowywanej lambdy},
ale nie przedłuża automatycznie czasu życia obiektów, do których lambda
odwołuje się przez referencje lub nieposiadające własności wskaźniki.

\begin{lstlisting}[caption={Callback utrzymujący obiekt przy życiu przez \texttt{shared\_ptr}; kompletny program.}]
#include <cassert>
#include <functional>
#include <memory>

int main() {
    auto state = std::make_shared<int>(10);
    std::function<int()> get_value = [state] {
        return *state;
    };
    state.reset();
    assert(get_value() == 10);
}
\end{lstlisting}

Trwałe przechwycenie \texttt{shared\_ptr} może z kolei utworzyć cykl
własności, np. gdy obiekt przechowuje callback, a callback przechwytuje
\texttt{shared\_ptr} do tego samego obiektu. Wówczas często lepiej
przechwycić \texttt{weak\_ptr} i uzyskać tymczasową własność przez
\texttt{lock()}.

\begin{lstlisting}[caption={Callback niewymuszający przedłużania życia obiektu; kompletny program.}]
#include <cassert>
#include <functional>
#include <memory>
#include <optional>

int main() {
    auto state = std::make_shared<int>(42);
    std::weak_ptr<int> weak = state;

    std::function<std::optional<int>()> read = [weak]()
        -> std::optional<int> {
        if (auto shared = weak.lock()) return *shared;
        return std::nullopt;
    };

    assert(read() == 42);
    state.reset();
    assert(!read().has_value());
}
\end{lstlisting}

\subsection{Szablon, wskaźnik do funkcji czy \texttt{std::function}?}

Jeżeli funkcja korzysta z obiektu wywoływalnego tylko podczas bieżącego
wywołania, często najprościej użyć szablonu. Pozwala to kompilatorowi
znać konkretny typ i daje okazję do \emph{inliningu}. Wskaźnik do funkcji
jest lekki, ale nie przechowa samodzielnie stanu przechwycenia lambdy.
\texttt{std::function} ułatwia natomiast przechowywanie i zastępowanie
heterogenicznych callbacków o wspólnej sygnaturze.

\begin{lstlisting}[caption={Wywołanie generyczne bez type erasure; kompletny program.}]
#include <cassert>
#include <functional>
#include <utility>

template<class F, class... Args>
decltype(auto) call_now(F&& f, Args&&... args) {
    return std::invoke(std::forward<F>(f),
                       std::forward<Args>(args)...);
}

int main() {
    int factor = 3;
    auto f = [factor](int x) { return factor * x; };
    assert(call_now(f, 7) == 21);
}
\end{lstlisting}

\section{Własny mechanizm type erasure: \texttt{Printable}}
\label{sec:cpp14-custom-erasure}

To rozszerzenie pokazuje, co może znajdować się pod gotowym
opakowaniem wymazującym typ. Przy pierwszej lekturze wystarczy
rozumieć, kiedy użyć \texttt{variant}, \texttt{any} lub
\texttt{function}; własną implementację można przeczytać później.

Pokażemy teraz, jak samodzielnie zbudować typ o stałym interfejsie,
który może przechowywać różne wartości drukowalne do strumienia. Celowo
nie wykorzystamy dziedziczenia typów użytkownika po wspólnej klasie
bazowej. Użytkownik powinien móc przekazać \texttt{int},
\texttt{std::string} albo własny typ, dla którego działa
\texttt{operator<<}. Wymagamy przy tym \emph{semantyki wartości}:
kopiowanie \texttt{Printable} ma kopiować przechowywany obiekt.

\paragraph{Idea implementacji.}
Wewnątrz opakowania znajdują się:
\begin{enumerate}
\item abstrakcyjny interfejs \texttt{Concept} z operacjami \texttt{print}
      i \texttt{clone};
\item szablonowy model \texttt{Model<T>}, który przechowuje konkretną
      wartość i realizuje interfejs;
\item wskaźnik \texttt{unique\_ptr<Concept>}, ukrywający konkretny
      typ w miejscu użycia \texttt{Printable}.
\end{enumerate}

\begin{lstlisting}[caption={Własne kopiowalne type erasure (C++20); kompletny program.}]
#include <cassert>
#include <concepts>
#include <memory>
#include <ostream>
#include <sstream>
#include <string>
#include <type_traits>
#include <utility>

class Printable {
    struct Concept {
        virtual ~Concept() = default;
        virtual void print(std::ostream&) const = 0;
        virtual std::unique_ptr<Concept> clone() const = 0;
    };

    template<class T>
    struct Model final : Concept {
        T value;
        template<class U>
        explicit Model(U&& x) : value(std::forward<U>(x)) {}
        void print(std::ostream& os) const override {
            os << value;
        }
        std::unique_ptr<Concept> clone() const override {
            return std::make_unique<Model<T>>(value);
        }
    };

    std::unique_ptr<Concept> self_;

public:
    Printable() = default;

    template<class T>
      requires (!std::same_as<std::remove_cvref_t<T>, Printable> &&
                std::copy_constructible<std::decay_t<T>> &&
                requires(std::ostream& os,
                         const std::decay_t<T>& value) {
                    { os << value } -> std::same_as<std::ostream&>;
                })
    Printable(T&& value)
        : self_(std::make_unique<Model<std::decay_t<T>>>(
              std::forward<T>(value))) {}

    Printable(const Printable& other)
        : self_(other.self_ ? other.self_->clone() : nullptr) {}

    Printable(Printable&&) noexcept = default;

    Printable& operator=(const Printable& other) {
        if (this != &other) {
            Printable tmp(other); // najpierw bezpieczna kopia
            swap(tmp);            // potem nie-rzucajaca zamiana
        }
        return *this;
    }

    Printable& operator=(Printable&&) noexcept = default;

    void swap(Printable& other) noexcept {
        self_.swap(other.self_);
    }

    void print(std::ostream& os) const {
        if (self_) self_->print(os);
        else os << "<empty>";
    }
};

int main() {
    Printable a = 17;
    Printable b = std::string("hello");
    Printable c = b; // gleboka kopia wartosci
    b = 7.5;

    std::ostringstream stream;
    a.print(stream);
    c.print(stream);
    assert(stream.str() == "17hello");
}
\end{lstlisting}

\subsection{Analiza własności i kosztów}

Ten interfejs jest \emph{otwarty}: każdy nowy typ spełniający wymagania
może zostać opakowany bez modyfikowania klasy \texttt{Printable}.
Użytkownik nie musi dziedziczyć po klasie bazowej, co jest ważną zaletą
w porównaniu ze zwykłym polimorfizmem opartym na \texttt{Base}.
\texttt{Model<T>} jest instancjonowany w czasie kompilacji, natomiast
\texttt{print()} jest wywoływane dynamicznie przez ukryty interfejs.
Kopiowanie przez \texttt{clone()} zachowuje rzeczywisty typ wartości.

Wersja edukacyjna wykorzystuje osobną alokację dynamiczną przy
przechowywaniu niepustej wartości i kolejną alokację przy głębokim
kopiowaniu. Nie wykorzystuje \emph{small-buffer optimization}, które
mogłoby przechowywać niewielkie obiekty bezpośrednio wewnątrz opakowania.
W produkcyjnym type erasure trzeba dodatkowo przemyśleć \emph{wyrównanie},
przenoszenie, bezpieczne niszczenie, usuwanie typów niekopiowalnych,
kwalifikatory \texttt{const}, obsługę wyjątków i stabilność ABI.

\paragraph{Gwarancja silna przypisania kopiującego.}
Tworzenie obiektu \texttt{tmp(other)} może zakończyć się wyjątkiem.
Wówczas obiekt docelowy pozostaje niezmieniony. Jeśli kopia się uda,
\texttt{swap(tmp)} jest operacją \texttt{noexcept}; stara zawartość
zostanie zniszczona razem z \texttt{tmp}. Jest to klasyczny idiom
\emph{copy-and-swap}. Alternatywą jest niekopiowalne opakowanie,
które usuwa operacje kopiowania i przechowuje tylko obiekty przenoszalne.

\section{Type erasure a dziedziczenie i polimorfizm statyczny}
\label{sec:cpp14-comparison}

\subsection{Wspólna baza klas}

Gdy typy zostały zaprojektowane jako część wspólnej hierarchii,
\texttt{std::unique\_ptr<Base>} pozwala przechowywać obiekty o różnych
typach dynamicznych i wywoływać ich operacje wirtualne. Implementacje
muszą dziedziczyć po \texttt{Base}, a tworzenie kopii polimorficznej
wymaga własnego protokołu (np. \texttt{clone()}). Z drugiej strony
\texttt{unique\_ptr} jasno wyraża własność i nie kopiuje obiektów
wskutek zwykłego kopiowania uchwytu (ponieważ kopiowanie jest zabronione).

\begin{lstlisting}[caption={Polimorfizm przez bazę i \texttt{unique\_ptr}; kompletny program.}]
#include <cassert>
#include <memory>
#include <string>
#include <vector>

struct Shape {
    virtual ~Shape() = default;
    virtual int area() const = 0;
};
struct Square final : Shape {
    int side;
    explicit Square(int s) : side(s) {}
    int area() const override { return side * side; }
};
struct Rectangle final : Shape {
    int width, height;
    Rectangle(int w, int h) : width(w), height(h) {}
    int area() const override { return width * height; }
};

int main() {
    std::vector<std::unique_ptr<Shape>> shapes;
    shapes.push_back(std::make_unique<Square>(3));
    shapes.push_back(std::make_unique<Rectangle>(2, 5));
    assert(shapes[0]->area() + shapes[1]->area() == 19);
}
\end{lstlisting}

\subsection{Kiedy nie wymazywać typu?}

Jeżeli algorytm potrzebuje konkretnego typu wyłącznie podczas jednego
wywołania, parametr szablonowy z \texttt{concept} jest często prostszy
i pozwala unikać alokacji. Type erasure ma sens przede wszystkim na
granicy interfejsów, gdy wartości o różnych typach muszą być
\emph{przechowywane w jednej strukturze}, przekazywane przez stabilny
interfejs albo zastępowane w trakcie działania programu.

\begin{center}
\small
\begin{tabularx}{\textwidth}{@{}lXX@{}}
\toprule
\textbf{Rozwiązanie} & \textbf{Zalety} & \textbf{Ograniczenia} \\
\midrule
\texttt{variant<Ts...>} & Typy znane statycznie, sprawdzanie przypadków, lokalne przechowywanie wartości & Zamknięty katalog typów, złożoność odwiedzania i rozmiar największej alternatywy \\
\texttt{any} & Dowolna kopiowalna wartość, wspólny typ pojemnika & Trzeba odzyskać dokładny typ; brak wspólnego interfejsu operacji \\
\texttt{function} & Jedna sygnatura dla różnych callbacków & Kopiowalność (C++20), pośrednictwo, możliwe alokacje \\
Własne type erasure & Otwarty katalog modeli bez wspólnej bazy typów klienta & Złożona implementacja, zarządzanie pamięcią i wyjątkami \\
\texttt{unique\_ptr<Base>} & Jawna hierarchia, dziedziczone operacje, kontrola własności & Wymaga wspólnej bazy; operacje wirtualne \\
Szablon z \texttt{concept} & Typ znany kompilatorowi, możliwość optymalizacji & Wiele instancjacji; zwykle nie przechowa różnych typów w jednej kolekcji \\
\bottomrule
\end{tabularx}
\end{center}

\section{Wydajność: układ pamięci i koszty dynamicznej dyspozycji}
\label{sec:cpp14-performance}

Nie ma uniwersalnej odpowiedzi na pytanie, który mechanizm jest
najszybszy. Poniżej wymieniamy dominujące koszty; faktyczne zachowanie
zależy od rozmiaru typów, optymalizacji kompilatora, struktury programu
i rozkładu typów w danych.

\begin{itemize}
\item \texttt{optional<T>} dodaje znacznik obecności; dostęp do zawartości
      może wymagać sprawdzenia warunku.
\item \texttt{variant<Ts...>} przechowuje wystarczającą ilość miejsca
      na największy typ i indeks aktywnej alternatywy; odwiedzający
      może być implementowany przez instrukcję wyboru lub tablicę
      funkcji, a szczegóły zależą od kompilatora.
\item \texttt{any} przechowuje metadane typu oraz wartość, potencjalnie
      w osobno alokowanym buforze.
\item \texttt{function} może wymagać pośredniego wywołania i przechowania
      stanu lambdy; czasami da się uniknąć alokacji.
\item Własna implementacja z \texttt{unique\_ptr<Concept>} wykonuje
      alokację dynamiczną i pośrednie wywołanie wirtualne.
\end{itemize}

Przede wszystkim trzeba mierzyć \emph{rzeczywisty} przypadek użycia.
Jednorazowy koszt konstruowania callbacku może być nieistotny, gdy jego
wykonanie trwa długo, ale może dominować w pętli składającej się z
milionów krótkich operacji.

\subsection{Test bez mylenia czasu kompilacji z wykonaniem}

Rzetelne porównanie wymaga rozgrzania kodu, wykorzystania odpowiedniego
zestawu danych, uwzględnienia optymalizacji i sprawdzenia, czy kompilator
nie usunął mierzonej pracy. Warto rozdzielać koszt \emph{tworzenia}
opakowania od kosztu \emph{wywołania} i analizować nie tylko średni
czas, ale też rozkład wyników oraz liczbę alokacji. Benchmarku
porównującego dwa mechanizmy nie traktujemy jako dowodu ich przewagi
we wszystkich zastosowaniach.

\section{Granice bezpieczeństwa: wyjątki, czas życia i współbieżność}

\paragraph{Wyjątki.}
\texttt{optional::value()}, \texttt{std::get} dla wariantu, \texttt{any\_cast}
i wywołanie pustego \texttt{std::function} mogą zgłaszać odpowiednie
wyjątki. Konstrukcja zawartości i przypisanie także mogą rzucać,
ponieważ wykonują operacje typu przechowywanego. Należy rozróżniać
\emph{brak wartości} (normalny stan \texttt{optional}) od
\emph{niepowodzenia operacji}, któremu warto nadać osobną reprezentację.

\paragraph{Czas życia.}
\texttt{variant}, \texttt{optional} i \texttt{any} posiadają własne
przechowywane wartości. Można jednak umieścić w nich wskaźniki,
\texttt{reference\_wrapper} albo inne nieposiadające własności uchwyty;
wtedy odpowiedzialność za czas życia wskazywanych danych pozostaje po
stronie programisty. To samo dotyczy referencji przechwyconych przez
lambdę umieszczoną w \texttt{std::function}.

\paragraph{Wątki.}
Żaden z opisanych typów nie czyni dowolnych współbieżnych modyfikacji
tego samego obiektu bezpiecznymi. Przykładowo równoczesny zapis do
jednego \texttt{variant} i odczyt jego aktywnej zawartości bez
synchronizacji mogą powodować wyścig danych. Gdy callback jest
uruchamiany równolegle, trzeba osobno określić zasady synchronizacji
stanu przechwyconego przez lambdę.

\section{Nowości po C++20 i świadomy wybór standardu}

W całym rozdziale przykłady są zgodne z C++20. W nowszych standardach
pojawiają się udogodnienia, które warto znać, ale których \emph{nie wolno}
zakładać przy kompilacji z \texttt{-std=c++20}:
\begin{itemize}
\item \textbf{C++23:} \texttt{std::expected<T,E>} --- wynik albo
      jawna wartość błędu;
\item \textbf{C++23:} \texttt{std::move\_only\_function} --- opakowanie
      dla obiektów wywoływalnych, które nie muszą być kopiowalne;
\item \textbf{C++23:} operacje monadyczne dla \texttt{std::optional}:
      \texttt{and\_then}, \texttt{transform} i \texttt{or\_else};
\item \textbf{C++26:} \texttt{std::function\_ref} --- niewłaścicielski
      widok na obiekt wywoływalny, którego czas życia musi zostać
      zapewniony przez wywołującego.
\end{itemize}

Dostępność funkcji bibliotecznych w konkretnym kompilatorze i jego
bibliotece standardowej może być opóźniona w stosunku do daty
opublikowania standardu.

\section{Zasady projektowania i typowe błędy}
\label{sec:cpp14-guidelines}

\begin{enumerate}
\item Wybieraj \texttt{optional<T>}, gdy brak wyniku jest normalny i
      nie potrzebujesz szczegółowego typu błędu.
\item Wybieraj \texttt{variant}, gdy zestaw przypadków jest znany
      i zamknięty; obsługuj je przez \texttt{visit}, jeśli pozwala to
      uniknąć łańcucha ręcznych kontroli typu.
\item Stosuj \texttt{any} tylko tam, gdzie świadomie rezygnujesz
      ze statycznie określonego zestawu typów i masz jasny protokół
      odzyskiwania ich wartości.
\item Używaj \texttt{std::function} do \emph{przechowywanych},
      kopiowalnych callbacków; dla funkcji wykonywanej natychmiast
      rozważ parametr szablonowy.
\item Nie zakładaj, że \texttt{std::function} zawsze alokuje ani
      że \texttt{std::variant} zawsze jest szybszy. Zmierz przypadek
      zastosowania.
\item Nie pozostawiaj referencji lub wskaźników do obiektów, których
      czas życia zakończy się przed użyciem przechowywanego callbacku.
\item Unikaj \texttt{std::get} i \texttt{any\_cast} w sytuacjach,
      w których nie masz gwarancji typu; do sprawdzania wykorzystuj
      \texttt{get\_if} i wskaźnikowy \texttt{any\_cast}.
\item Nie używaj \texttt{std::any} do bezpośredniego opakowania
      typów przenoszalnych, ale niekopiowalnych, np. \texttt{unique\_ptr}.
\item Rozdziel decyzję o \emph{otwartym zbiorze implementacji} od
      decyzji o \emph{dynamicznym czasie życia i alokacji}.
\end{enumerate}

\section*{Podsumowanie rozdziału}
\addcontentsline{toc}{section}{Podsumowanie rozdziału 10}

\texttt{optional} modeluje brak lub obecność wartości, a
\texttt{variant} modeluje zamkniętą sumę przypadków. \texttt{any}
ukrywa typ kopiowalnej wartości, nie narzucając jej operacji.
\texttt{function} ukrywa typ obiektu wywoływalnego przy zachowaniu
ustalonej sygnatury. Własne type erasure pozwala opakować otwartą
rodzinę obcych typów bez wymagania wspólnej klasy bazowej, ale wymaga
starannego określenia semantyki własności, kopiowania i wyjątków.
Dobry projekt zaczyna się więc od określenia kontraktu: \emph{które
przypadki istnieją, kto posiada obiekty, jakie operacje udostępniamy
oraz czy zbiór typów ma pozostać otwarty}.

\section*{Dokumentacja i literatura}
\addcontentsline{toc}{section}{Dokumentacja i literatura rozdziału 10}
\begin{itemize}
\item \url{https://en.cppreference.com/w/cpp/utility/optional} --- \texttt{std::optional}.
\item \url{https://en.cppreference.com/w/cpp/utility/variant} --- \texttt{std::variant}.
\item \url{https://en.cppreference.com/w/cpp/utility/any} --- \texttt{std::any}.
\item \href{https://en.cppreference.com/w/cpp/utility/functional/function}{cppreference: \texttt{std::function}} --- dokumentacja opakowania wywołań.
\item Bjarne Stroustrup, \emph{A Tour of C++}, wyd. 3, Addison--Wesley, 2022.
\item Klaus Iglberger, \emph{C++ Software Design}, O'Reilly, 2022.
\end{itemize}

% Kompletne programy sa oznaczone w podpisach listingow.

\chapter{Wymazywanie typów i typy sumaryczne}
\label{chap:cpp14-type-erasure}

\section{Problem: jak reprezentować różne rodzaje wartości?}

W programowaniu generycznym często chcemy posługiwać się różnymi typami
przez wspólny interfejs. C++ oferuje kilka mechanizmów realizacji tego celu,
ale mechanizmy te nie są zamienne.
Wyobraźmy sobie aplikację z trzema poleceniami: zapisaniem pliku,
dodaniem dwóch liczb i zakończeniem pracy. Gdy znamy wszystkie trzy
rodzaje poleceń, możemy je wymienić w \texttt{std::variant}. Jeśli
zewnętrzne moduły mają dodawać nowe polecenia bez zmiany tej listy,
potrzebujemy otwartego interfejsu, np. klasy bazowej z funkcją wirtualną.
Wybór mechanizmu zaczynamy więc od pytań:
\begin{enumerate}
\item Czy zestaw dozwolonych typów jest \emph{zamknięty} (ustalony z góry),
      czy \emph{otwarty} (może zostać rozszerzony o nowy typ)?
\item Czy chcemy przechować \emph{wartość}, czy jedynie referencję/uchwyt
      do istniejącego obiektu?
\item Czy wywoływane operacje mają być znane już podczas kompilacji,
      czy wybierane dynamicznie w czasie działania programu?
\item Czy dopuszczamy alokację dynamiczną, kopiowanie oraz koszt
      pośredniego wywołania funkcji?
\end{enumerate}

\emph{Typ sumaryczny} (ang. \emph{sum type}) reprezentuje alternatywę:
wartość należy do jednego z kilku przypadków. W C++17 taką rolę pełni
\texttt{std::variant<Ts...>}; \texttt{std::optional<T>} opisuje wariant
,,brak wartości albo wartość \texttt{T}''. Natomiast \emph{type erasure}
(\emph{wymazywanie typu}) pozwala ukryć konkretny typ obiektu za
interfejsem nieparametryzowanym przez ten typ. Standardowe przykłady
to \texttt{std::any} i \texttt{std::function}.

\begin{center}
\small
\begin{tabularx}{\textwidth}{@{}lXX@{}}
\toprule
\textbf{Mechanizm} & \textbf{Informacja o typie} & \textbf{Zasadnicze zastosowanie} \\
\midrule
\texttt{optional<T>} & Jedna możliwość \texttt{T} lub brak & Wynik opcjonalny, wartość nieobecna \\
\texttt{variant<Ts...>} & Zamknięty zbiór alternatyw & Zdarzenia, komendy, model stanu \\
\texttt{any} & Typ ustalany w czasie wykonania & Heterogeniczne wartości bez wspólnego interfejsu operacji \\
\texttt{function<R(Args...)>} & Dowolny zgodny obiekt wywoływalny & Przechowywany callback, funkcja jako wartość \\
\texttt{unique\_ptr<Base>} & Otwarta hierarchia klas & Polimorfizm dynamiczny i tożsamość obiektu \\
Szablon/\texttt{concept} & Typ znany przy instancjowaniu & Polimorfizm statyczny, specjalizacja kodu \\
\bottomrule
\end{tabularx}
\end{center}

Najprostsza ściąga to: \texttt{optional<int>} oznacza liczbę albo jej
brak, \texttt{variant<int, std::string>} jedną z dwóch znanych
postaci wartości, a \texttt{any} wartość o typie, którego autor
pojemnika nie musi znać. \texttt{function<void()>} przechowuje
różne działania, które można wywołać w taki sam sposób. W kolejnych
sekcjach zobaczymy, co z tych wyborów wynika dla sprawdzania typów,
własności danych i kosztu wykonania.

\section{\texttt{std::optional<T>}: wartość, której może nie być}

\begin{center}\small
\begin{tabularx}{\textwidth}{@{}p{0.33\textwidth}X@{}}
\toprule
Typ lub operacja & Co reprezentuje \\
\midrule
\code{std::optional<T>} & Wartość typu \code{T} albo jej brak; nie opisuje przyczyny błędu. \\
\code{has_value()} & Sprawdza, czy \code{optional} przechowuje wartość. \\
\code{value()} & Zwraca wartość lub zgłasza \code{std::bad_optional_access}, gdy jej brak. \\
\code{value_or(fallback)} & Zwraca wartość albo podany zamiennik, konwertowany do \code{T}. \\
\code{std::variant<Ts...>} & Dokładnie jedną z ustalonych alternatyw typów. \\
\code{std::any} & Wartość typu ustalonego dopiero w czasie wykonania; przy odczycie trzeba sprawdzić typ. \\
\bottomrule
\end{tabularx}
\end{center}
\label{sec:cpp14-optional}

Typ \texttt{std::optional<T>} (C++17, nagłówek \texttt{<optional>})
przechowuje bezpośrednio obiekt typu \texttt{T}, jeżeli jest on obecny,
lub reprezentuje stan pusty. Jest to pełnoprawna \emph{wartość}, a nie
wskaźnik. W szczególności pusty \texttt{optional} nie oznacza wskaźnika
zerowego, a niepusty nie wymaga automatycznie osobnej alokacji dynamicznej.

\begin{lstlisting}[caption={Najważniejsze operacje optional (C++17); kompletny program.}]
#include <cassert>
#include <optional>
#include <string>

int main() {
    std::optional<int> x;         // stan pusty
    assert(!x.has_value());
    assert(x == std::nullopt);

    x = 42;
    assert(x.has_value());
    assert(*x == 42);
    assert(x.value() == 42);
    assert(x.value_or(-1) == 42);

    x.reset();
    assert(x.value_or(-1) == -1);
    x.emplace(7);
    assert(*x == 7);

    std::optional<std::string> name{std::in_place, 3, 'a'};
    assert(*name == "aaa");
}
\end{lstlisting}

Dostęp przez \texttt{*x} i \texttt{x->...} wymaga obecności wartości:
\emph{użycie pustego obiektu w taki sposób jest niezdefiniowanym
zachowaniem}. Metoda \texttt{value()} wykonuje sprawdzenie i dla stanu
pustego zgłasza \texttt{std::bad\_optional\_access}. Warto pamiętać,
że \texttt{value\_or(default)} przyjmuje wartość domyślną jako argument
funkcji: wyrażenie tworzące argument jest obliczane również wtedy,
gdy \texttt{optional} zawiera wartość. Może ponadto następować kopiowanie
lub przeniesienie przechowywanego obiektu.

\subsection{Praktyczny przykład: parsowanie liczby}

Funkcja parsująca dane wejściowe może zwrócić \texttt{optional}, jeśli
interesuje nas tylko informacja ,,udało się/nie udało się''.
\texttt{std::from\_chars} unika wyjątków i nie wymaga tworzenia dodatkowego
łańcucha zakończonego zerem. W przykładzie wymagamy, aby \emph{cały}
łańcuch był prawidłową liczbą.

\begin{lstlisting}[caption={Parsowanie int do optional (C++17); kompletny program.}]
#include <cassert>
#include <charconv>
#include <optional>
#include <string_view>
#include <system_error>

std::optional<int> parse_int(std::string_view text) {
    int result = 0;
    auto [ptr, ec] = std::from_chars(
        text.data(), text.data() + text.size(), result);
    if (ec != std::errc{} || ptr != text.data() + text.size())
        return std::nullopt;
    return result;
}

int main() {
    assert(parse_int("123") == 123);
    assert(parse_int("123x") == std::nullopt);
    assert(parse_int("") == std::nullopt);
    assert(parse_int("99999999999999999999") == std::nullopt);
}
\end{lstlisting}

\subsection{Własność, referencje i koszt}

\texttt{optional<T>} posiada swoją wartość i zwykle zajmuje więcej pamięci
niż sam \texttt{T}, ponieważ musi przechowywać informację, czy wartość
istnieje. Rozmiar zależy od typu i implementacji. Standard C++20
nie udostępnia \texttt{std::optional<T\&>}. Jeśli wynik ma wskazywać na
istniejący obiekt, można użyć wskaźnika nieposiadającego własności albo
\texttt{std::optional<std::reference\_wrapper<T>>}. W obu przypadkach
programista odpowiada za czas życia wskazywanego obiektu.

\begin{lstlisting}[caption={Opcjonalna referencja: brak przenoszenia własności; fragment.}]
#include <functional>
#include <optional>
#include <string>

std::optional<std::reference_wrapper<const std::string>>
find_name(bool exists, const std::string& name) {
    if (exists) return std::cref(name);
    return std::nullopt;
}
// Zwracana referencja jest poprawna tylko gdy `name` nadal istnieje.
\end{lstlisting}

Gdy brak wartości jest sytuacją normalną, \texttt{optional} najczęściej
czytelniej wyraża intencję niż wyjątkowe zakończenie funkcji. Jeżeli
konieczna jest informacja \emph{dlaczego} nie uzyskano wyniku, należy
zaprojektować osobny typ błędu, użyć \texttt{variant} lub (w C++23)
\texttt{std::expected<T,E>}.

\section{\texttt{std::variant<Ts...>}: zamknięta alternatywa typów}
\label{sec:cpp14-variant}

\texttt{std::variant<Ts...>} (C++17, nagłówek \texttt{<variant>})
przechowuje \emph{jedną aktywną alternatywę} z listy typów \texttt{Ts...}.
Nie jest odpowiednikiem \texttt{std::tuple<Ts...>}: krotka zawiera
\emph{wszystkie} składowe, a wariant dokładnie jeden z dopuszczonych
przypadków (poza szczególnym stanem \texttt{valueless\_by\_exception()}).

\begin{lstlisting}[caption={Tworzenie i badanie aktywnej alternatywy (C++17); kompletny program.}]
#include <cassert>
#include <string>
#include <variant>

int main() {
    using Value = std::variant<int, std::string>;
    Value x = 7;
    assert(x.index() == 0);
    assert(std::holds_alternative<int>(x));
    assert(std::get<int>(x) == 7);

    x = std::string("hello");
    assert(x.index() == 1);
    assert(std::get_if<int>(&x) == nullptr);
    assert(std::get_if<std::string>(&x) != nullptr);

    x.emplace<0>(19);
    assert(std::get<0>(x) == 19);
}
\end{lstlisting}

Niepoprawne \texttt{std::get<T>(x)} (gdy aktywny jest inny typ)
zgłasza \texttt{std::bad\_variant\_access}. Alternatywa
\texttt{std::get\_if<T>(\&x)} zwraca \texttt{nullptr}, a zatem nadaje się
do sprawdzania bez wyjątku. \texttt{std::get<T>} i
\texttt{std::holds\_alternative<T>} wymagają, aby \texttt{T} występował
\emph{dokładnie raz} w liście alternatyw; dla powtarzających się typów
używamy indeksów. Nie należy także zapominać o możliwym niepowodzeniu
operacji konwersji przy konstruowaniu wariantu.

\subsection{Domyślna alternatywa i \texttt{std::monostate}}

Domyślnie konstruowany wariant próbuje skonstruować pierwszą alternatywę.
Jeżeli chcemy jawnie reprezentować stan pusty, można umieścić na początku
\texttt{std::monostate}; zapewnia on także domyślną konstruowalność
wariantu, gdy inny pierwszy typ nie jest domyślnie konstruowalny.

\begin{lstlisting}[caption={Jawny stan początkowy monostate; kompletny program.}]
#include <cassert>
#include <string>
#include <variant>

int main() {
    using State = std::variant<std::monostate, int, std::string>;
    State s;
    assert(std::holds_alternative<std::monostate>(s));
    s = std::string("ready");
    assert(std::get<std::string>(s) == "ready");
}
\end{lstlisting}

\subsection{Odwiedzanie wariantu: \texttt{std::visit}}

Samo wielokrotne sprawdzanie \texttt{holds\_alternative} szybko
prowadzi do rozbudowanych instrukcji \texttt{if}. Mechanizm
\texttt{std::visit(visitor, variant)} wywołuje odwiedzającego z aktywną
wartością odpowiedniego typu. Gdy odwiedzający jest lambdą generyczną,
kompilator instancjuje jej operator wywołania dla różnych alternatyw.

\begin{lstlisting}[caption={Jeden odwiedzający generyczny; kompletny program.}]
#include <cassert>
#include <cstddef>
#include <string>
#include <type_traits>
#include <variant>

int main() {
    std::variant<int, double, std::string> v = "hi";

    const std::size_t size = std::visit([](const auto& x) {
        using T = std::decay_t<decltype(x)>;
        if constexpr (std::is_same_v<T, std::string>)
            return x.size();
        else
            return sizeof(x);
    }, v);

    assert(size == 2);
}
\end{lstlisting}

Do osobnej obsługi poszczególnych typów wygodny jest wzorzec
\texttt{overloaded}, znany również z rozdziału o lambdach.
Wykorzystuje on dziedziczenie po kilku obiektach funkcyjnych oraz
ujawnienie wszystkich operatorów wywołania za pomocą \texttt{using}.

\begin{lstlisting}[caption={overloaded oraz std::visit (C++17); kompletny program.}]
#include <cassert>
#include <cstddef>
#include <string>
#include <variant>

template<class... Fs>
struct overloaded : Fs... {
    using Fs::operator()...;
};
template<class... Fs> overloaded(Fs...) -> overloaded<Fs...>;

int main() {
    using Data = std::variant<int, double, std::string>;
    Data data = std::string("abc");

    auto length = std::visit(overloaded{
        [](int) { return std::size_t{1}; },
        [](double) { return std::size_t{1}; },
        [](const std::string& s) { return s.size(); }
    }, data);
    assert(length == 3);
}
\end{lstlisting}

W przypadku odwiedzania zwracającego wartość wymagany jest zgodny typ
wyniku wszystkich możliwych wywołań (w praktyce należy zapewnić
identyczny typ i kategorię wartości zwracanych wyrażeń, z wyjątkiem
szczególnych reguł dla \texttt{void}). Przy wielu wariantach
\texttt{std::visit} uwzględnia \emph{kombinacje} alternatyw. To daje
czytelny kod przy małej liczbie przypadków, ale może zwiększyć liczbę
instancjacji szablonów.

\begin{lstlisting}[caption={Odwiedzenie dwóch wariantów naraz; kompletny program.}]
#include <cassert>
#include <variant>

int main() {
    std::variant<int, double> a = 3;
    std::variant<int, double> b = 2.5;
    const double sum = std::visit(
        [](const auto& x, const auto& y) -> double {
            return static_cast<double>(x) + static_cast<double>(y);
        }, a, b);
    assert(sum == 5.5);
}
\end{lstlisting}

\subsection{Modyfikacja stanu i kwalifikacja \texttt{const}}

\texttt{std::visit} respektuje typ i kategorię wartości przekazanego
wariantu. Odwiedzając obiekt przez \texttt{variant\&}, możemy modyfikować
aktywną alternatywę; przez \texttt{const variant\&} uzyskujemy
referencję stałą. Przy przekazaniu \texttt{std::move(v)} odwiedzający
może otrzymać referencję rvalue i przenieść przechowywaną wartość.
Przeniesienie nie musi zmieniać aktywnego indeksu wariantu, ale zawartość
przeniesionego obiektu może zostać zmieniona zgodnie z semantyką typu.

\begin{lstlisting}[caption={Modyfikacja aktywnej alternatywy; kompletny program.}]
#include <cassert>
#include <variant>

int main() {
    std::variant<int, double> v = 10;
    std::visit([](auto& x) { x += 1; }, v);
    assert(std::get<int>(v) == 11);
}
\end{lstlisting}

\subsection{Stan \texttt{valueless\_by\_exception}}

Wyjątek podczas przypisania lub tworzenia nowej wartości może
pozostawić \texttt{variant} bez aktywnej alternatywy. Ten stan
wykrywa metoda \texttt{valueless\_by\_exception()}.
Wtedy \texttt{index()} zwraca \texttt{std::variant\_npos}.
Wywołanie \texttt{std::visit} zgłasza \texttt{std::bad\_variant\_access}.
Nie jest to zwykła ,,pusta'' alternatywa. Jeśli jest potrzebna,
można użyć \texttt{std::monostate}.

\paragraph{Ważne ograniczenia.}
\texttt{variant} nie przechowuje typów referencyjnych ani \texttt{void};
nie można także dodać dowolnego nowego typu bez zmiany listy alternatyw.
Duży wariant może zajmować pamięć odpowiadającą największej alternatywie
plus informację o aktywnym przypadku i ewentualne wyrównanie; dokładny
rozmiar nie jest stałą wielkością określoną przez standard.

\section{Przykład projektowy: interpreter prostych poleceń}
\label{sec:cpp14-commands}

Rozpatrzmy miniaturowy model poleceń. Każdy rodzaj polecenia ma inne pola,
ale wszystkie należą do zamkniętej rodziny. \texttt{variant} oddziela
\emph{dane} od \emph{operacji}, a \texttt{visit} pozwala jednoznacznie
obsłużyć każdy przypadek.

\begin{lstlisting}[caption={Modelowanie poleceń za pomocą variant; kompletny program.}]
#include <cassert>
#include <string>
#include <variant>
#include <vector>

struct Add  { int a, b; };
struct Save { std::string filename; };
struct Quit {};
using Command = std::variant<Add, Save, Quit>;

struct Execute {
    int& total;
    bool& stopped;

    void operator()(const Add& c) const {
        total += c.a + c.b;
    }
    void operator()(const Save& c) const {
        if (c.filename.empty()) stopped = true;
    }
    void operator()(const Quit&) const {
        stopped = true;
    }
};

int main() {
    std::vector<Command> commands{
        Add{2, 3}, Save{"result.txt"}, Add{4, 5}, Quit{}
    };
    int total = 0;
    bool stopped = false;
    for (const auto& command : commands) {
        if (stopped) break;
        std::visit(Execute{total, stopped}, command);
    }
    assert(total == 14);
    assert(stopped);
}
\end{lstlisting}

Dodanie nowego typu polecenia, np. \texttt{Delete}, wymaga rozszerzenia
\texttt{Command} i obsługi nowego przypadku. W dobrze skonstruowanym
odwiedzającym brakująca obsługa będzie błędem kompilacji. To zaleta
typów sumarycznych przy modelowaniu \emph{zamkniętego} katalogu zdarzeń.

\section{\texttt{std::any}: wartość o nieznanym typie}
\label{sec:cpp14-any}

\texttt{std::any} (C++17, nagłówek \texttt{<any>}) może przechowywać
obiekt dowolnego typu spełniającego wymaganie kopiowalności i pozwala
później odzyskać dokładny typ przez \texttt{std::any\_cast}. To klasyczna
implementacja \emph{type erasure}: użytkownik obiektu \texttt{any}
nie zna szablonowego typu przechowywanej wartości.

\begin{lstlisting}[caption={std::any i bezpieczne wydobycie typu (C++17); kompletny program.}]
#include <any>
#include <cassert>
#include <string>

int main() {
    std::any value = 17;
    assert(value.has_value());
    assert(std::any_cast<int>(value) == 17);

    if (auto ptr = std::any_cast<int>(&value)) {
        *ptr += 1;
    }
    assert(std::any_cast<int>(value) == 18);

    value.emplace<std::string>("text");
    assert(std::any_cast<int>(&value) == nullptr);
    assert(std::any_cast<const std::string&>(value) == "text");

    value.reset();
    assert(!value.has_value());
}
\end{lstlisting}

\texttt{std::any\_cast<T>(value)} zgłasza \texttt{std::bad\_any\_cast},
jeśli aktywny typ nie odpowiada żądanemu. Wariant wskaźnikowy
\texttt{std::any\_cast<T>(\&value)} zwraca \texttt{nullptr}; jest zatem
wygodny do ostrożnego sprawdzania typu. \texttt{type()} udostępnia
\texttt{std::type\_info} opisujący aktualnie przechowywany typ,
a dla pustego obiektu zwraca \texttt{typeid(void)}.

\begin{lstlisting}[caption={\texttt{any\_cast} nie jest konwersją numeryczną; kompletny program.}]
#include <any>
#include <cassert>
#include <typeinfo>

int main() {
    std::any value = 42;
    assert(value.type() == typeid(int));
    assert(std::any_cast<double>(&value) == nullptr);
    assert(std::any_cast<int>(&value) != nullptr);
}
\end{lstlisting}

\texttt{any} przechowuje \emph{wartość}, nie udostępnia jednak uniwersalnego
operatora \texttt{+}, porównania czy drukowania tej wartości. Aby wykonać
operację, trzeba znać typ albo dostarczyć odrębny mechanizm dyspozycji.
Ponadto przechowywany typ musi być kopiowalny: \texttt{std::any}
nie nadaje się do bezpośredniego przechowywania \texttt{std::unique\_ptr<T>}.
Niektóre implementacje wykorzystują \emph{small-object optimization},
ale standard nie obiecuje uniknięcia alokacji dla konkretnego małego typu.

\paragraph{Kiedy używać?}
\texttt{std::any} ma sens m.in. w luźno typowanych zasobnikach kontekstu,
gdy autor zasobnika nie zna listy możliwych typów, a odbiorca umie
jawnie odtworzyć typ właściwego wpisu. Jeśli rodzaje wartości są znane,
\texttt{variant} daje z reguły silniejszą kontrolę statyczną i czytelniejszą
obsługę przypadków. Unikaj \texttt{any} jako sposobu na ukrycie
źle określonych kontraktów interfejsu.

\section{\texttt{std::function}: wymazywanie typu obiektu wywoływalnego}
\label{sec:cpp14-function}

\texttt{std::function<R(Args...)>} (C++11, nagłówek
\texttt{<functional>}) przechowuje kopiowalny obiekt wywoływalny zgodny
z podaną sygnaturą. Może to być wskaźnik do funkcji, funktor,
wyrażenie lambda, a także obiekt związany z metodą składową. Konkretny
typ obiektu zostaje ukryty, ale interfejs jego wywołania jest znany:
\texttt{R(Args...)}.

\begin{lstlisting}[caption={Jeden typ std::function dla różnych obiektów wywoływalnych; kompletny program.}]
#include <cassert>
#include <functional>

int times_two(int x) { return 2 * x; }
struct PlusThree {
    int operator()(int x) const { return x + 3; }
};

int main() {
    std::function<int(int)> f = times_two;
    assert(f(4) == 8);

    f = PlusThree{};
    assert(f(4) == 7);

    int delta = 5;
    f = [delta](int x) { return x + delta; };
    assert(f(4) == 9);
}
\end{lstlisting}

\subsection{Pusty obiekt i koszt abstrakcji}

Domyślnie skonstruowany \texttt{std::function} jest pusty. Jego
\texttt{operator bool()} informuje, czy istnieje obiekt wywoływalny;
próba wywołania pustego obiektu zgłasza \texttt{std::bad\_function\_call}.
W porównaniu z bezpośrednim wywołaniem szablonu \texttt{std::function}
może wymagać pośredniego wywołania i dodatkowego przechowywania stanu,
a czasami także alokacji dynamicznej. Nie jest jednak poprawne twierdzenie,
że każdy \texttt{std::function} \emph{zawsze} wykonuje alokację.

\begin{lstlisting}[caption={Pusty std::function i jawne sprawdzanie; kompletny program.}]
#include <cassert>
#include <functional>

int main() {
    std::function<void()> callback;
    assert(!callback);
    callback = [] {};
    assert(static_cast<bool>(callback));
    callback();
    callback = nullptr;
    assert(!callback);
}
\end{lstlisting}

\subsection{Ograniczenie kopiowalności i czas życia przechwyceń}

W C++20 \texttt{std::function} wymaga kopiowalności typu przechowywanego
obiektu wywoływalnego. Dlatego lambda zawierająca \emph{move-only}
\texttt{std::unique\_ptr} nie może być bezpośrednio umieszczona w
\texttt{std::function}. W C++23 pojawił się
\texttt{std::move\_only\_function} dla takiego zastosowania.

Przechwytywanie przez referencję wymaga ponadto, aby obiekt odniesienia
pozostawał żywy podczas każdego późniejszego wywołania. Sam
\texttt{std::function} kontroluje czas życia \emph{przechowywanej lambdy},
ale nie przedłuża automatycznie czasu życia obiektów, do których lambda
odwołuje się przez referencje lub nieposiadające własności wskaźniki.

\begin{lstlisting}[caption={Callback utrzymujący obiekt przy życiu przez \texttt{shared\_ptr}; kompletny program.}]
#include <cassert>
#include <functional>
#include <memory>

int main() {
    auto state = std::make_shared<int>(10);
    std::function<int()> get_value = [state] {
        return *state;
    };
    state.reset();
    assert(get_value() == 10);
}
\end{lstlisting}

Trwałe przechwycenie \texttt{shared\_ptr} może z kolei utworzyć cykl
własności, np. gdy obiekt przechowuje callback, a callback przechwytuje
\texttt{shared\_ptr} do tego samego obiektu. Wówczas często lepiej
przechwycić \texttt{weak\_ptr} i uzyskać tymczasową własność przez
\texttt{lock()}.

\begin{lstlisting}[caption={Callback niewymuszający przedłużania życia obiektu; kompletny program.}]
#include <cassert>
#include <functional>
#include <memory>
#include <optional>

int main() {
    auto state = std::make_shared<int>(42);
    std::weak_ptr<int> weak = state;

    std::function<std::optional<int>()> read = [weak]()
        -> std::optional<int> {
        if (auto shared = weak.lock()) return *shared;
        return std::nullopt;
    };

    assert(read() == 42);
    state.reset();
    assert(!read().has_value());
}
\end{lstlisting}

\subsection{Szablon, wskaźnik do funkcji czy \texttt{std::function}?}

Jeżeli funkcja korzysta z obiektu wywoływalnego tylko podczas bieżącego
wywołania, często najprościej użyć szablonu. Pozwala to kompilatorowi
znać konkretny typ i daje okazję do \emph{inliningu}. Wskaźnik do funkcji
jest lekki, ale nie przechowa samodzielnie stanu przechwycenia lambdy.
\texttt{std::function} ułatwia natomiast przechowywanie i zastępowanie
heterogenicznych callbacków o wspólnej sygnaturze.

\begin{lstlisting}[caption={Wywołanie generyczne bez type erasure; kompletny program.}]
#include <cassert>
#include <functional>
#include <utility>

template<class F, class... Args>
decltype(auto) call_now(F&& f, Args&&... args) {
    return std::invoke(std::forward<F>(f),
                       std::forward<Args>(args)...);
}

int main() {
    int factor = 3;
    auto f = [factor](int x) { return factor * x; };
    assert(call_now(f, 7) == 21);
}
\end{lstlisting}

\section{Własny mechanizm type erasure: \texttt{Printable}}
\label{sec:cpp14-custom-erasure}

To rozszerzenie pokazuje, co może znajdować się pod gotowym
opakowaniem wymazującym typ. Przy pierwszej lekturze wystarczy
rozumieć, kiedy użyć \texttt{variant}, \texttt{any} lub
\texttt{function}; własną implementację można przeczytać później.

Pokażemy teraz, jak samodzielnie zbudować typ o stałym interfejsie,
który może przechowywać różne wartości drukowalne do strumienia. Celowo
nie wykorzystamy dziedziczenia typów użytkownika po wspólnej klasie
bazowej. Użytkownik powinien móc przekazać \texttt{int},
\texttt{std::string} albo własny typ, dla którego działa
\texttt{operator<<}. Wymagamy przy tym \emph{semantyki wartości}:
kopiowanie \texttt{Printable} ma kopiować przechowywany obiekt.

\paragraph{Idea implementacji.}
Wewnątrz opakowania znajdują się:
\begin{enumerate}
\item abstrakcyjny interfejs \texttt{Concept} z operacjami \texttt{print}
      i \texttt{clone};
\item szablonowy model \texttt{Model<T>}, który przechowuje konkretną
      wartość i realizuje interfejs;
\item wskaźnik \texttt{unique\_ptr<Concept>}, ukrywający konkretny
      typ w miejscu użycia \texttt{Printable}.
\end{enumerate}

\begin{lstlisting}[caption={Własne kopiowalne type erasure (C++20); kompletny program.}]
#include <cassert>
#include <concepts>
#include <memory>
#include <ostream>
#include <sstream>
#include <string>
#include <type_traits>
#include <utility>

class Printable {
    struct Concept {
        virtual ~Concept() = default;
        virtual void print(std::ostream&) const = 0;
        virtual std::unique_ptr<Concept> clone() const = 0;
    };

    template<class T>
    struct Model final : Concept {
        T value;
        template<class U>
        explicit Model(U&& x) : value(std::forward<U>(x)) {}
        void print(std::ostream& os) const override {
            os << value;
        }
        std::unique_ptr<Concept> clone() const override {
            return std::make_unique<Model<T>>(value);
        }
    };

    std::unique_ptr<Concept> self_;

public:
    Printable() = default;

    template<class T>
      requires (!std::same_as<std::remove_cvref_t<T>, Printable> &&
                std::copy_constructible<std::decay_t<T>> &&
                requires(std::ostream& os,
                         const std::decay_t<T>& value) {
                    { os << value } -> std::same_as<std::ostream&>;
                })
    Printable(T&& value)
        : self_(std::make_unique<Model<std::decay_t<T>>>(
              std::forward<T>(value))) {}

    Printable(const Printable& other)
        : self_(other.self_ ? other.self_->clone() : nullptr) {}

    Printable(Printable&&) noexcept = default;

    Printable& operator=(const Printable& other) {
        if (this != &other) {
            Printable tmp(other); // najpierw bezpieczna kopia
            swap(tmp);            // potem nie-rzucajaca zamiana
        }
        return *this;
    }

    Printable& operator=(Printable&&) noexcept = default;

    void swap(Printable& other) noexcept {
        self_.swap(other.self_);
    }

    void print(std::ostream& os) const {
        if (self_) self_->print(os);
        else os << "<empty>";
    }
};

int main() {
    Printable a = 17;
    Printable b = std::string("hello");
    Printable c = b; // gleboka kopia wartosci
    b = 7.5;

    std::ostringstream stream;
    a.print(stream);
    c.print(stream);
    assert(stream.str() == "17hello");
}
\end{lstlisting}

\subsection{Analiza własności i kosztów}

Ten interfejs jest \emph{otwarty}: każdy nowy typ spełniający wymagania
może zostać opakowany bez modyfikowania klasy \texttt{Printable}.
Użytkownik nie musi dziedziczyć po klasie bazowej, co jest ważną zaletą
w porównaniu ze zwykłym polimorfizmem opartym na \texttt{Base}.
\texttt{Model<T>} jest instancjonowany w czasie kompilacji, natomiast
\texttt{print()} jest wywoływane dynamicznie przez ukryty interfejs.
Kopiowanie przez \texttt{clone()} zachowuje rzeczywisty typ wartości.

Wersja edukacyjna wykorzystuje osobną alokację dynamiczną przy
przechowywaniu niepustej wartości i kolejną alokację przy głębokim
kopiowaniu. Nie wykorzystuje \emph{small-buffer optimization}, które
mogłoby przechowywać niewielkie obiekty bezpośrednio wewnątrz opakowania.
W produkcyjnym type erasure trzeba dodatkowo przemyśleć \emph{wyrównanie},
przenoszenie, bezpieczne niszczenie, usuwanie typów niekopiowalnych,
kwalifikatory \texttt{const}, obsługę wyjątków i stabilność ABI.

\paragraph{Gwarancja silna przypisania kopiującego.}
Tworzenie obiektu \texttt{tmp(other)} może zakończyć się wyjątkiem.
Wówczas obiekt docelowy pozostaje niezmieniony. Jeśli kopia się uda,
\texttt{swap(tmp)} jest operacją \texttt{noexcept}; stara zawartość
zostanie zniszczona razem z \texttt{tmp}. Jest to klasyczny idiom
\emph{copy-and-swap}. Alternatywą jest niekopiowalne opakowanie,
które usuwa operacje kopiowania i przechowuje tylko obiekty przenoszalne.

\section{Type erasure a dziedziczenie i polimorfizm statyczny}
\label{sec:cpp14-comparison}

\subsection{Wspólna baza klas}

Gdy typy zostały zaprojektowane jako część wspólnej hierarchii,
\texttt{std::unique\_ptr<Base>} pozwala przechowywać obiekty o różnych
typach dynamicznych i wywoływać ich operacje wirtualne. Implementacje
muszą dziedziczyć po \texttt{Base}, a tworzenie kopii polimorficznej
wymaga własnego protokołu (np. \texttt{clone()}). Z drugiej strony
\texttt{unique\_ptr} jasno wyraża własność i nie kopiuje obiektów
wskutek zwykłego kopiowania uchwytu (ponieważ kopiowanie jest zabronione).

\begin{lstlisting}[caption={Polimorfizm przez bazę i \texttt{unique\_ptr}; kompletny program.}]
#include <cassert>
#include <memory>
#include <string>
#include <vector>

struct Shape {
    virtual ~Shape() = default;
    virtual int area() const = 0;
};
struct Square final : Shape {
    int side;
    explicit Square(int s) : side(s) {}
    int area() const override { return side * side; }
};
struct Rectangle final : Shape {
    int width, height;
    Rectangle(int w, int h) : width(w), height(h) {}
    int area() const override { return width * height; }
};

int main() {
    std::vector<std::unique_ptr<Shape>> shapes;
    shapes.push_back(std::make_unique<Square>(3));
    shapes.push_back(std::make_unique<Rectangle>(2, 5));
    assert(shapes[0]->area() + shapes[1]->area() == 19);
}
\end{lstlisting}

\subsection{Kiedy nie wymazywać typu?}

Jeżeli algorytm potrzebuje konkretnego typu wyłącznie podczas jednego
wywołania, parametr szablonowy z \texttt{concept} jest często prostszy
i pozwala unikać alokacji. Type erasure ma sens przede wszystkim na
granicy interfejsów, gdy wartości o różnych typach muszą być
\emph{przechowywane w jednej strukturze}, przekazywane przez stabilny
interfejs albo zastępowane w trakcie działania programu.

\begin{center}
\small
\begin{tabularx}{\textwidth}{@{}lXX@{}}
\toprule
\textbf{Rozwiązanie} & \textbf{Zalety} & \textbf{Ograniczenia} \\
\midrule
\texttt{variant<Ts...>} & Typy znane statycznie, sprawdzanie przypadków, lokalne przechowywanie wartości & Zamknięty katalog typów, złożoność odwiedzania i rozmiar największej alternatywy \\
\texttt{any} & Dowolna kopiowalna wartość, wspólny typ pojemnika & Trzeba odzyskać dokładny typ; brak wspólnego interfejsu operacji \\
\texttt{function} & Jedna sygnatura dla różnych callbacków & Kopiowalność (C++20), pośrednictwo, możliwe alokacje \\
Własne type erasure & Otwarty katalog modeli bez wspólnej bazy typów klienta & Złożona implementacja, zarządzanie pamięcią i wyjątkami \\
\texttt{unique\_ptr<Base>} & Jawna hierarchia, dziedziczone operacje, kontrola własności & Wymaga wspólnej bazy; operacje wirtualne \\
Szablon z \texttt{concept} & Typ znany kompilatorowi, możliwość optymalizacji & Wiele instancjacji; zwykle nie przechowa różnych typów w jednej kolekcji \\
\bottomrule
\end{tabularx}
\end{center}

\section{Wydajność: układ pamięci i koszty dynamicznej dyspozycji}
\label{sec:cpp14-performance}

Nie ma uniwersalnej odpowiedzi na pytanie, który mechanizm jest
najszybszy. Poniżej wymieniamy dominujące koszty; faktyczne zachowanie
zależy od rozmiaru typów, optymalizacji kompilatora, struktury programu
i rozkładu typów w danych.

\begin{itemize}
\item \texttt{optional<T>} dodaje znacznik obecności; dostęp do zawartości
      może wymagać sprawdzenia warunku.
\item \texttt{variant<Ts...>} przechowuje wystarczającą ilość miejsca
      na największy typ i indeks aktywnej alternatywy; odwiedzający
      może być implementowany przez instrukcję wyboru lub tablicę
      funkcji, a szczegóły zależą od kompilatora.
\item \texttt{any} przechowuje metadane typu oraz wartość, potencjalnie
      w osobno alokowanym buforze.
\item \texttt{function} może wymagać pośredniego wywołania i przechowania
      stanu lambdy; czasami da się uniknąć alokacji.
\item Własna implementacja z \texttt{unique\_ptr<Concept>} wykonuje
      alokację dynamiczną i pośrednie wywołanie wirtualne.
\end{itemize}

Przede wszystkim trzeba mierzyć \emph{rzeczywisty} przypadek użycia.
Jednorazowy koszt konstruowania callbacku może być nieistotny, gdy jego
wykonanie trwa długo, ale może dominować w pętli składającej się z
milionów krótkich operacji.

\subsection{Test bez mylenia czasu kompilacji z wykonaniem}

Rzetelne porównanie wymaga rozgrzania kodu, wykorzystania odpowiedniego
zestawu danych, uwzględnienia optymalizacji i sprawdzenia, czy kompilator
nie usunął mierzonej pracy. Warto rozdzielać koszt \emph{tworzenia}
opakowania od kosztu \emph{wywołania} i analizować nie tylko średni
czas, ale też rozkład wyników oraz liczbę alokacji. Benchmarku
porównującego dwa mechanizmy nie traktujemy jako dowodu ich przewagi
we wszystkich zastosowaniach.

\section{Granice bezpieczeństwa: wyjątki, czas życia i współbieżność}

\paragraph{Wyjątki.}
\texttt{optional::value()}, \texttt{std::get} dla wariantu, \texttt{any\_cast}
i wywołanie pustego \texttt{std::function} mogą zgłaszać odpowiednie
wyjątki. Konstrukcja zawartości i przypisanie także mogą rzucać,
ponieważ wykonują operacje typu przechowywanego. Należy rozróżniać
\emph{brak wartości} (normalny stan \texttt{optional}) od
\emph{niepowodzenia operacji}, któremu warto nadać osobną reprezentację.

\paragraph{Czas życia.}
\texttt{variant}, \texttt{optional} i \texttt{any} posiadają własne
przechowywane wartości. Można jednak umieścić w nich wskaźniki,
\texttt{reference\_wrapper} albo inne nieposiadające własności uchwyty;
wtedy odpowiedzialność za czas życia wskazywanych danych pozostaje po
stronie programisty. To samo dotyczy referencji przechwyconych przez
lambdę umieszczoną w \texttt{std::function}.

\paragraph{Wątki.}
Żaden z opisanych typów nie czyni dowolnych współbieżnych modyfikacji
tego samego obiektu bezpiecznymi. Przykładowo równoczesny zapis do
jednego \texttt{variant} i odczyt jego aktywnej zawartości bez
synchronizacji mogą powodować wyścig danych. Gdy callback jest
uruchamiany równolegle, trzeba osobno określić zasady synchronizacji
stanu przechwyconego przez lambdę.

\section{Nowości po C++20 i świadomy wybór standardu}

W całym rozdziale przykłady są zgodne z C++20. W nowszych standardach
pojawiają się udogodnienia, które warto znać, ale których \emph{nie wolno}
zakładać przy kompilacji z \texttt{-std=c++20}:
\begin{itemize}
\item \textbf{C++23:} \texttt{std::expected<T,E>} --- wynik albo
      jawna wartość błędu;
\item \textbf{C++23:} \texttt{std::move\_only\_function} --- opakowanie
      dla obiektów wywoływalnych, które nie muszą być kopiowalne;
\item \textbf{C++23:} operacje monadyczne dla \texttt{std::optional}:
      \texttt{and\_then}, \texttt{transform} i \texttt{or\_else};
\item \textbf{C++26:} \texttt{std::function\_ref} --- niewłaścicielski
      widok na obiekt wywoływalny, którego czas życia musi zostać
      zapewniony przez wywołującego.
\end{itemize}

Dostępność funkcji bibliotecznych w konkretnym kompilatorze i jego
bibliotece standardowej może być opóźniona w stosunku do daty
opublikowania standardu.

\section{Zasady projektowania i typowe błędy}
\label{sec:cpp14-guidelines}

\begin{enumerate}
\item Wybieraj \texttt{optional<T>}, gdy brak wyniku jest normalny i
      nie potrzebujesz szczegółowego typu błędu.
\item Wybieraj \texttt{variant}, gdy zestaw przypadków jest znany
      i zamknięty; obsługuj je przez \texttt{visit}, jeśli pozwala to
      uniknąć łańcucha ręcznych kontroli typu.
\item Stosuj \texttt{any} tylko tam, gdzie świadomie rezygnujesz
      ze statycznie określonego zestawu typów i masz jasny protokół
      odzyskiwania ich wartości.
\item Używaj \texttt{std::function} do \emph{przechowywanych},
      kopiowalnych callbacków; dla funkcji wykonywanej natychmiast
      rozważ parametr szablonowy.
\item Nie zakładaj, że \texttt{std::function} zawsze alokuje ani
      że \texttt{std::variant} zawsze jest szybszy. Zmierz przypadek
      zastosowania.
\item Nie pozostawiaj referencji lub wskaźników do obiektów, których
      czas życia zakończy się przed użyciem przechowywanego callbacku.
\item Unikaj \texttt{std::get} i \texttt{any\_cast} w sytuacjach,
      w których nie masz gwarancji typu; do sprawdzania wykorzystuj
      \texttt{get\_if} i wskaźnikowy \texttt{any\_cast}.
\item Nie używaj \texttt{std::any} do bezpośredniego opakowania
      typów przenoszalnych, ale niekopiowalnych, np. \texttt{unique\_ptr}.
\item Rozdziel decyzję o \emph{otwartym zbiorze implementacji} od
      decyzji o \emph{dynamicznym czasie życia i alokacji}.
\end{enumerate}

\section*{Podsumowanie rozdziału}
\addcontentsline{toc}{section}{Podsumowanie rozdziału 10}

\texttt{optional} modeluje brak lub obecność wartości, a
\texttt{variant} modeluje zamkniętą sumę przypadków. \texttt{any}
ukrywa typ kopiowalnej wartości, nie narzucając jej operacji.
\texttt{function} ukrywa typ obiektu wywoływalnego przy zachowaniu
ustalonej sygnatury. Własne type erasure pozwala opakować otwartą
rodzinę obcych typów bez wymagania wspólnej klasy bazowej, ale wymaga
starannego określenia semantyki własności, kopiowania i wyjątków.
Dobry projekt zaczyna się więc od określenia kontraktu: \emph{które
przypadki istnieją, kto posiada obiekty, jakie operacje udostępniamy
oraz czy zbiór typów ma pozostać otwarty}.

\section*{Dokumentacja i literatura}
\addcontentsline{toc}{section}{Dokumentacja i literatura rozdziału 10}
\begin{itemize}
\item \url{https://en.cppreference.com/w/cpp/utility/optional} --- \texttt{std::optional}.
\item \url{https://en.cppreference.com/w/cpp/utility/variant} --- \texttt{std::variant}.
\item \url{https://en.cppreference.com/w/cpp/utility/any} --- \texttt{std::any}.
\item \href{https://en.cppreference.com/w/cpp/utility/functional/function}{cppreference: \texttt{std::function}} --- dokumentacja opakowania wywołań.
\item Bjarne Stroustrup, \emph{A Tour of C++}, wyd. 3, Addison--Wesley, 2022.
\item Klaus Iglberger, \emph{C++ Software Design}, O'Reilly, 2022.
\end{itemize}

% Kompletne programy sa oznaczone w podpisach listingow.

\chapter{Wymazywanie typów i typy sumaryczne}
\label{chap:cpp14-type-erasure}

\section{Problem: jak reprezentować różne rodzaje wartości?}

W programowaniu generycznym często chcemy posługiwać się różnymi typami
przez wspólny interfejs. C++ oferuje kilka mechanizmów realizacji tego celu,
ale mechanizmy te nie są zamienne.
Wyobraźmy sobie aplikację z trzema poleceniami: zapisaniem pliku,
dodaniem dwóch liczb i zakończeniem pracy. Gdy znamy wszystkie trzy
rodzaje poleceń, możemy je wymienić w \texttt{std::variant}. Jeśli
zewnętrzne moduły mają dodawać nowe polecenia bez zmiany tej listy,
potrzebujemy otwartego interfejsu, np. klasy bazowej z funkcją wirtualną.
Wybór mechanizmu zaczynamy więc od pytań:
\begin{enumerate}
\item Czy zestaw dozwolonych typów jest \emph{zamknięty} (ustalony z góry),
      czy \emph{otwarty} (może zostać rozszerzony o nowy typ)?
\item Czy chcemy przechować \emph{wartość}, czy jedynie referencję/uchwyt
      do istniejącego obiektu?
\item Czy wywoływane operacje mają być znane już podczas kompilacji,
      czy wybierane dynamicznie w czasie działania programu?
\item Czy dopuszczamy alokację dynamiczną, kopiowanie oraz koszt
      pośredniego wywołania funkcji?
\end{enumerate}

\emph{Typ sumaryczny} (ang. \emph{sum type}) reprezentuje alternatywę:
wartość należy do jednego z kilku przypadków. W C++17 taką rolę pełni
\texttt{std::variant<Ts...>}; \texttt{std::optional<T>} opisuje wariant
,,brak wartości albo wartość \texttt{T}''. Natomiast \emph{type erasure}
(\emph{wymazywanie typu}) pozwala ukryć konkretny typ obiektu za
interfejsem nieparametryzowanym przez ten typ. Standardowe przykłady
to \texttt{std::any} i \texttt{std::function}.

\begin{center}
\small
\begin{tabularx}{\textwidth}{@{}lXX@{}}
\toprule
\textbf{Mechanizm} & \textbf{Informacja o typie} & \textbf{Zasadnicze zastosowanie} \\
\midrule
\texttt{optional<T>} & Jedna możliwość \texttt{T} lub brak & Wynik opcjonalny, wartość nieobecna \\
\texttt{variant<Ts...>} & Zamknięty zbiór alternatyw & Zdarzenia, komendy, model stanu \\
\texttt{any} & Typ ustalany w czasie wykonania & Heterogeniczne wartości bez wspólnego interfejsu operacji \\
\texttt{function<R(Args...)>} & Dowolny zgodny obiekt wywoływalny & Przechowywany callback, funkcja jako wartość \\
\texttt{unique\_ptr<Base>} & Otwarta hierarchia klas & Polimorfizm dynamiczny i tożsamość obiektu \\
Szablon/\texttt{concept} & Typ znany przy instancjowaniu & Polimorfizm statyczny, specjalizacja kodu \\
\bottomrule
\end{tabularx}
\end{center}

Najprostsza ściąga to: \texttt{optional<int>} oznacza liczbę albo jej
brak, \texttt{variant<int, std::string>} jedną z dwóch znanych
postaci wartości, a \texttt{any} wartość o typie, którego autor
pojemnika nie musi znać. \texttt{function<void()>} przechowuje
różne działania, które można wywołać w taki sam sposób. W kolejnych
sekcjach zobaczymy, co z tych wyborów wynika dla sprawdzania typów,
własności danych i kosztu wykonania.

\section{\texttt{std::optional<T>}: wartość, której może nie być}

\begin{center}\small
\begin{tabularx}{\textwidth}{@{}p{0.33\textwidth}X@{}}
\toprule
Typ lub operacja & Co reprezentuje \\
\midrule
\code{std::optional<T>} & Wartość typu \code{T} albo jej brak; nie opisuje przyczyny błędu. \\
\code{has_value()} & Sprawdza, czy \code{optional} przechowuje wartość. \\
\code{value()} & Zwraca wartość lub zgłasza \code{std::bad_optional_access}, gdy jej brak. \\
\code{value_or(fallback)} & Zwraca wartość albo podany zamiennik, konwertowany do \code{T}. \\
\code{std::variant<Ts...>} & Dokładnie jedną z ustalonych alternatyw typów. \\
\code{std::any} & Wartość typu ustalonego dopiero w czasie wykonania; przy odczycie trzeba sprawdzić typ. \\
\bottomrule
\end{tabularx}
\end{center}
\label{sec:cpp14-optional}

Typ \texttt{std::optional<T>} (C++17, nagłówek \texttt{<optional>})
przechowuje bezpośrednio obiekt typu \texttt{T}, jeżeli jest on obecny,
lub reprezentuje stan pusty. Jest to pełnoprawna \emph{wartość}, a nie
wskaźnik. W szczególności pusty \texttt{optional} nie oznacza wskaźnika
zerowego, a niepusty nie wymaga automatycznie osobnej alokacji dynamicznej.

\begin{lstlisting}[caption={Najważniejsze operacje optional (C++17); kompletny program.}]
#include <cassert>
#include <optional>
#include <string>

int main() {
    std::optional<int> x;         // stan pusty
    assert(!x.has_value());
    assert(x == std::nullopt);

    x = 42;
    assert(x.has_value());
    assert(*x == 42);
    assert(x.value() == 42);
    assert(x.value_or(-1) == 42);

    x.reset();
    assert(x.value_or(-1) == -1);
    x.emplace(7);
    assert(*x == 7);

    std::optional<std::string> name{std::in_place, 3, 'a'};
    assert(*name == "aaa");
}
\end{lstlisting}

Dostęp przez \texttt{*x} i \texttt{x->...} wymaga obecności wartości:
\emph{użycie pustego obiektu w taki sposób jest niezdefiniowanym
zachowaniem}. Metoda \texttt{value()} wykonuje sprawdzenie i dla stanu
pustego zgłasza \texttt{std::bad\_optional\_access}. Warto pamiętać,
że \texttt{value\_or(default)} przyjmuje wartość domyślną jako argument
funkcji: wyrażenie tworzące argument jest obliczane również wtedy,
gdy \texttt{optional} zawiera wartość. Może ponadto następować kopiowanie
lub przeniesienie przechowywanego obiektu.

\subsection{Praktyczny przykład: parsowanie liczby}

Funkcja parsująca dane wejściowe może zwrócić \texttt{optional}, jeśli
interesuje nas tylko informacja ,,udało się/nie udało się''.
\texttt{std::from\_chars} unika wyjątków i nie wymaga tworzenia dodatkowego
łańcucha zakończonego zerem. W przykładzie wymagamy, aby \emph{cały}
łańcuch był prawidłową liczbą.

\begin{lstlisting}[caption={Parsowanie int do optional (C++17); kompletny program.}]
#include <cassert>
#include <charconv>
#include <optional>
#include <string_view>
#include <system_error>

std::optional<int> parse_int(std::string_view text) {
    int result = 0;
    auto [ptr, ec] = std::from_chars(
        text.data(), text.data() + text.size(), result);
    if (ec != std::errc{} || ptr != text.data() + text.size())
        return std::nullopt;
    return result;
}

int main() {
    assert(parse_int("123") == 123);
    assert(parse_int("123x") == std::nullopt);
    assert(parse_int("") == std::nullopt);
    assert(parse_int("99999999999999999999") == std::nullopt);
}
\end{lstlisting}

\subsection{Własność, referencje i koszt}

\texttt{optional<T>} posiada swoją wartość i zwykle zajmuje więcej pamięci
niż sam \texttt{T}, ponieważ musi przechowywać informację, czy wartość
istnieje. Rozmiar zależy od typu i implementacji. Standard C++20
nie udostępnia \texttt{std::optional<T\&>}. Jeśli wynik ma wskazywać na
istniejący obiekt, można użyć wskaźnika nieposiadającego własności albo
\texttt{std::optional<std::reference\_wrapper<T>>}. W obu przypadkach
programista odpowiada za czas życia wskazywanego obiektu.

\begin{lstlisting}[caption={Opcjonalna referencja: brak przenoszenia własności; fragment.}]
#include <functional>
#include <optional>
#include <string>

std::optional<std::reference_wrapper<const std::string>>
find_name(bool exists, const std::string& name) {
    if (exists) return std::cref(name);
    return std::nullopt;
}
// Zwracana referencja jest poprawna tylko gdy `name` nadal istnieje.
\end{lstlisting}

Gdy brak wartości jest sytuacją normalną, \texttt{optional} najczęściej
czytelniej wyraża intencję niż wyjątkowe zakończenie funkcji. Jeżeli
konieczna jest informacja \emph{dlaczego} nie uzyskano wyniku, należy
zaprojektować osobny typ błędu, użyć \texttt{variant} lub (w C++23)
\texttt{std::expected<T,E>}.

\section{\texttt{std::variant<Ts...>}: zamknięta alternatywa typów}
\label{sec:cpp14-variant}

\texttt{std::variant<Ts...>} (C++17, nagłówek \texttt{<variant>})
przechowuje \emph{jedną aktywną alternatywę} z listy typów \texttt{Ts...}.
Nie jest odpowiednikiem \texttt{std::tuple<Ts...>}: krotka zawiera
\emph{wszystkie} składowe, a wariant dokładnie jeden z dopuszczonych
przypadków (poza szczególnym stanem \texttt{valueless\_by\_exception()}).

\begin{lstlisting}[caption={Tworzenie i badanie aktywnej alternatywy (C++17); kompletny program.}]
#include <cassert>
#include <string>
#include <variant>

int main() {
    using Value = std::variant<int, std::string>;
    Value x = 7;
    assert(x.index() == 0);
    assert(std::holds_alternative<int>(x));
    assert(std::get<int>(x) == 7);

    x = std::string("hello");
    assert(x.index() == 1);
    assert(std::get_if<int>(&x) == nullptr);
    assert(std::get_if<std::string>(&x) != nullptr);

    x.emplace<0>(19);
    assert(std::get<0>(x) == 19);
}
\end{lstlisting}

Niepoprawne \texttt{std::get<T>(x)} (gdy aktywny jest inny typ)
zgłasza \texttt{std::bad\_variant\_access}. Alternatywa
\texttt{std::get\_if<T>(\&x)} zwraca \texttt{nullptr}, a zatem nadaje się
do sprawdzania bez wyjątku. \texttt{std::get<T>} i
\texttt{std::holds\_alternative<T>} wymagają, aby \texttt{T} występował
\emph{dokładnie raz} w liście alternatyw; dla powtarzających się typów
używamy indeksów. Nie należy także zapominać o możliwym niepowodzeniu
operacji konwersji przy konstruowaniu wariantu.

\subsection{Domyślna alternatywa i \texttt{std::monostate}}

Domyślnie konstruowany wariant próbuje skonstruować pierwszą alternatywę.
Jeżeli chcemy jawnie reprezentować stan pusty, można umieścić na początku
\texttt{std::monostate}; zapewnia on także domyślną konstruowalność
wariantu, gdy inny pierwszy typ nie jest domyślnie konstruowalny.

\begin{lstlisting}[caption={Jawny stan początkowy monostate; kompletny program.}]
#include <cassert>
#include <string>
#include <variant>

int main() {
    using State = std::variant<std::monostate, int, std::string>;
    State s;
    assert(std::holds_alternative<std::monostate>(s));
    s = std::string("ready");
    assert(std::get<std::string>(s) == "ready");
}
\end{lstlisting}

\subsection{Odwiedzanie wariantu: \texttt{std::visit}}

Samo wielokrotne sprawdzanie \texttt{holds\_alternative} szybko
prowadzi do rozbudowanych instrukcji \texttt{if}. Mechanizm
\texttt{std::visit(visitor, variant)} wywołuje odwiedzającego z aktywną
wartością odpowiedniego typu. Gdy odwiedzający jest lambdą generyczną,
kompilator instancjuje jej operator wywołania dla różnych alternatyw.

\begin{lstlisting}[caption={Jeden odwiedzający generyczny; kompletny program.}]
#include <cassert>
#include <cstddef>
#include <string>
#include <type_traits>
#include <variant>

int main() {
    std::variant<int, double, std::string> v = "hi";

    const std::size_t size = std::visit([](const auto& x) {
        using T = std::decay_t<decltype(x)>;
        if constexpr (std::is_same_v<T, std::string>)
            return x.size();
        else
            return sizeof(x);
    }, v);

    assert(size == 2);
}
\end{lstlisting}

Do osobnej obsługi poszczególnych typów wygodny jest wzorzec
\texttt{overloaded}, znany również z rozdziału o lambdach.
Wykorzystuje on dziedziczenie po kilku obiektach funkcyjnych oraz
ujawnienie wszystkich operatorów wywołania za pomocą \texttt{using}.

\begin{lstlisting}[caption={overloaded oraz std::visit (C++17); kompletny program.}]
#include <cassert>
#include <cstddef>
#include <string>
#include <variant>

template<class... Fs>
struct overloaded : Fs... {
    using Fs::operator()...;
};
template<class... Fs> overloaded(Fs...) -> overloaded<Fs...>;

int main() {
    using Data = std::variant<int, double, std::string>;
    Data data = std::string("abc");

    auto length = std::visit(overloaded{
        [](int) { return std::size_t{1}; },
        [](double) { return std::size_t{1}; },
        [](const std::string& s) { return s.size(); }
    }, data);
    assert(length == 3);
}
\end{lstlisting}

W przypadku odwiedzania zwracającego wartość wymagany jest zgodny typ
wyniku wszystkich możliwych wywołań (w praktyce należy zapewnić
identyczny typ i kategorię wartości zwracanych wyrażeń, z wyjątkiem
szczególnych reguł dla \texttt{void}). Przy wielu wariantach
\texttt{std::visit} uwzględnia \emph{kombinacje} alternatyw. To daje
czytelny kod przy małej liczbie przypadków, ale może zwiększyć liczbę
instancjacji szablonów.

\begin{lstlisting}[caption={Odwiedzenie dwóch wariantów naraz; kompletny program.}]
#include <cassert>
#include <variant>

int main() {
    std::variant<int, double> a = 3;
    std::variant<int, double> b = 2.5;
    const double sum = std::visit(
        [](const auto& x, const auto& y) -> double {
            return static_cast<double>(x) + static_cast<double>(y);
        }, a, b);
    assert(sum == 5.5);
}
\end{lstlisting}

\subsection{Modyfikacja stanu i kwalifikacja \texttt{const}}

\texttt{std::visit} respektuje typ i kategorię wartości przekazanego
wariantu. Odwiedzając obiekt przez \texttt{variant\&}, możemy modyfikować
aktywną alternatywę; przez \texttt{const variant\&} uzyskujemy
referencję stałą. Przy przekazaniu \texttt{std::move(v)} odwiedzający
może otrzymać referencję rvalue i przenieść przechowywaną wartość.
Przeniesienie nie musi zmieniać aktywnego indeksu wariantu, ale zawartość
przeniesionego obiektu może zostać zmieniona zgodnie z semantyką typu.

\begin{lstlisting}[caption={Modyfikacja aktywnej alternatywy; kompletny program.}]
#include <cassert>
#include <variant>

int main() {
    std::variant<int, double> v = 10;
    std::visit([](auto& x) { x += 1; }, v);
    assert(std::get<int>(v) == 11);
}
\end{lstlisting}

\subsection{Stan \texttt{valueless\_by\_exception}}

Wyjątek podczas przypisania lub tworzenia nowej wartości może
pozostawić \texttt{variant} bez aktywnej alternatywy. Ten stan
wykrywa metoda \texttt{valueless\_by\_exception()}.
Wtedy \texttt{index()} zwraca \texttt{std::variant\_npos}.
Wywołanie \texttt{std::visit} zgłasza \texttt{std::bad\_variant\_access}.
Nie jest to zwykła ,,pusta'' alternatywa. Jeśli jest potrzebna,
można użyć \texttt{std::monostate}.

\paragraph{Ważne ograniczenia.}
\texttt{variant} nie przechowuje typów referencyjnych ani \texttt{void};
nie można także dodać dowolnego nowego typu bez zmiany listy alternatyw.
Duży wariant może zajmować pamięć odpowiadającą największej alternatywie
plus informację o aktywnym przypadku i ewentualne wyrównanie; dokładny
rozmiar nie jest stałą wielkością określoną przez standard.

\section{Przykład projektowy: interpreter prostych poleceń}
\label{sec:cpp14-commands}

Rozpatrzmy miniaturowy model poleceń. Każdy rodzaj polecenia ma inne pola,
ale wszystkie należą do zamkniętej rodziny. \texttt{variant} oddziela
\emph{dane} od \emph{operacji}, a \texttt{visit} pozwala jednoznacznie
obsłużyć każdy przypadek.

\begin{lstlisting}[caption={Modelowanie poleceń za pomocą variant; kompletny program.}]
#include <cassert>
#include <string>
#include <variant>
#include <vector>

struct Add  { int a, b; };
struct Save { std::string filename; };
struct Quit {};
using Command = std::variant<Add, Save, Quit>;

struct Execute {
    int& total;
    bool& stopped;

    void operator()(const Add& c) const {
        total += c.a + c.b;
    }
    void operator()(const Save& c) const {
        if (c.filename.empty()) stopped = true;
    }
    void operator()(const Quit&) const {
        stopped = true;
    }
};

int main() {
    std::vector<Command> commands{
        Add{2, 3}, Save{"result.txt"}, Add{4, 5}, Quit{}
    };
    int total = 0;
    bool stopped = false;
    for (const auto& command : commands) {
        if (stopped) break;
        std::visit(Execute{total, stopped}, command);
    }
    assert(total == 14);
    assert(stopped);
}
\end{lstlisting}

Dodanie nowego typu polecenia, np. \texttt{Delete}, wymaga rozszerzenia
\texttt{Command} i obsługi nowego przypadku. W dobrze skonstruowanym
odwiedzającym brakująca obsługa będzie błędem kompilacji. To zaleta
typów sumarycznych przy modelowaniu \emph{zamkniętego} katalogu zdarzeń.

\section{\texttt{std::any}: wartość o nieznanym typie}
\label{sec:cpp14-any}

\texttt{std::any} (C++17, nagłówek \texttt{<any>}) może przechowywać
obiekt dowolnego typu spełniającego wymaganie kopiowalności i pozwala
później odzyskać dokładny typ przez \texttt{std::any\_cast}. To klasyczna
implementacja \emph{type erasure}: użytkownik obiektu \texttt{any}
nie zna szablonowego typu przechowywanej wartości.

\begin{lstlisting}[caption={std::any i bezpieczne wydobycie typu (C++17); kompletny program.}]
#include <any>
#include <cassert>
#include <string>

int main() {
    std::any value = 17;
    assert(value.has_value());
    assert(std::any_cast<int>(value) == 17);

    if (auto ptr = std::any_cast<int>(&value)) {
        *ptr += 1;
    }
    assert(std::any_cast<int>(value) == 18);

    value.emplace<std::string>("text");
    assert(std::any_cast<int>(&value) == nullptr);
    assert(std::any_cast<const std::string&>(value) == "text");

    value.reset();
    assert(!value.has_value());
}
\end{lstlisting}

\texttt{std::any\_cast<T>(value)} zgłasza \texttt{std::bad\_any\_cast},
jeśli aktywny typ nie odpowiada żądanemu. Wariant wskaźnikowy
\texttt{std::any\_cast<T>(\&value)} zwraca \texttt{nullptr}; jest zatem
wygodny do ostrożnego sprawdzania typu. \texttt{type()} udostępnia
\texttt{std::type\_info} opisujący aktualnie przechowywany typ,
a dla pustego obiektu zwraca \texttt{typeid(void)}.

\begin{lstlisting}[caption={\texttt{any\_cast} nie jest konwersją numeryczną; kompletny program.}]
#include <any>
#include <cassert>
#include <typeinfo>

int main() {
    std::any value = 42;
    assert(value.type() == typeid(int));
    assert(std::any_cast<double>(&value) == nullptr);
    assert(std::any_cast<int>(&value) != nullptr);
}
\end{lstlisting}

\texttt{any} przechowuje \emph{wartość}, nie udostępnia jednak uniwersalnego
operatora \texttt{+}, porównania czy drukowania tej wartości. Aby wykonać
operację, trzeba znać typ albo dostarczyć odrębny mechanizm dyspozycji.
Ponadto przechowywany typ musi być kopiowalny: \texttt{std::any}
nie nadaje się do bezpośredniego przechowywania \texttt{std::unique\_ptr<T>}.
Niektóre implementacje wykorzystują \emph{small-object optimization},
ale standard nie obiecuje uniknięcia alokacji dla konkretnego małego typu.

\paragraph{Kiedy używać?}
\texttt{std::any} ma sens m.in. w luźno typowanych zasobnikach kontekstu,
gdy autor zasobnika nie zna listy możliwych typów, a odbiorca umie
jawnie odtworzyć typ właściwego wpisu. Jeśli rodzaje wartości są znane,
\texttt{variant} daje z reguły silniejszą kontrolę statyczną i czytelniejszą
obsługę przypadków. Unikaj \texttt{any} jako sposobu na ukrycie
źle określonych kontraktów interfejsu.

\section{\texttt{std::function}: wymazywanie typu obiektu wywoływalnego}
\label{sec:cpp14-function}

\texttt{std::function<R(Args...)>} (C++11, nagłówek
\texttt{<functional>}) przechowuje kopiowalny obiekt wywoływalny zgodny
z podaną sygnaturą. Może to być wskaźnik do funkcji, funktor,
wyrażenie lambda, a także obiekt związany z metodą składową. Konkretny
typ obiektu zostaje ukryty, ale interfejs jego wywołania jest znany:
\texttt{R(Args...)}.

\begin{lstlisting}[caption={Jeden typ std::function dla różnych obiektów wywoływalnych; kompletny program.}]
#include <cassert>
#include <functional>

int times_two(int x) { return 2 * x; }
struct PlusThree {
    int operator()(int x) const { return x + 3; }
};

int main() {
    std::function<int(int)> f = times_two;
    assert(f(4) == 8);

    f = PlusThree{};
    assert(f(4) == 7);

    int delta = 5;
    f = [delta](int x) { return x + delta; };
    assert(f(4) == 9);
}
\end{lstlisting}

\subsection{Pusty obiekt i koszt abstrakcji}

Domyślnie skonstruowany \texttt{std::function} jest pusty. Jego
\texttt{operator bool()} informuje, czy istnieje obiekt wywoływalny;
próba wywołania pustego obiektu zgłasza \texttt{std::bad\_function\_call}.
W porównaniu z bezpośrednim wywołaniem szablonu \texttt{std::function}
może wymagać pośredniego wywołania i dodatkowego przechowywania stanu,
a czasami także alokacji dynamicznej. Nie jest jednak poprawne twierdzenie,
że każdy \texttt{std::function} \emph{zawsze} wykonuje alokację.

\begin{lstlisting}[caption={Pusty std::function i jawne sprawdzanie; kompletny program.}]
#include <cassert>
#include <functional>

int main() {
    std::function<void()> callback;
    assert(!callback);
    callback = [] {};
    assert(static_cast<bool>(callback));
    callback();
    callback = nullptr;
    assert(!callback);
}
\end{lstlisting}

\subsection{Ograniczenie kopiowalności i czas życia przechwyceń}

W C++20 \texttt{std::function} wymaga kopiowalności typu przechowywanego
obiektu wywoływalnego. Dlatego lambda zawierająca \emph{move-only}
\texttt{std::unique\_ptr} nie może być bezpośrednio umieszczona w
\texttt{std::function}. W C++23 pojawił się
\texttt{std::move\_only\_function} dla takiego zastosowania.

Przechwytywanie przez referencję wymaga ponadto, aby obiekt odniesienia
pozostawał żywy podczas każdego późniejszego wywołania. Sam
\texttt{std::function} kontroluje czas życia \emph{przechowywanej lambdy},
ale nie przedłuża automatycznie czasu życia obiektów, do których lambda
odwołuje się przez referencje lub nieposiadające własności wskaźniki.

\begin{lstlisting}[caption={Callback utrzymujący obiekt przy życiu przez \texttt{shared\_ptr}; kompletny program.}]
#include <cassert>
#include <functional>
#include <memory>

int main() {
    auto state = std::make_shared<int>(10);
    std::function<int()> get_value = [state] {
        return *state;
    };
    state.reset();
    assert(get_value() == 10);
}
\end{lstlisting}

Trwałe przechwycenie \texttt{shared\_ptr} może z kolei utworzyć cykl
własności, np. gdy obiekt przechowuje callback, a callback przechwytuje
\texttt{shared\_ptr} do tego samego obiektu. Wówczas często lepiej
przechwycić \texttt{weak\_ptr} i uzyskać tymczasową własność przez
\texttt{lock()}.

\begin{lstlisting}[caption={Callback niewymuszający przedłużania życia obiektu; kompletny program.}]
#include <cassert>
#include <functional>
#include <memory>
#include <optional>

int main() {
    auto state = std::make_shared<int>(42);
    std::weak_ptr<int> weak = state;

    std::function<std::optional<int>()> read = [weak]()
        -> std::optional<int> {
        if (auto shared = weak.lock()) return *shared;
        return std::nullopt;
    };

    assert(read() == 42);
    state.reset();
    assert(!read().has_value());
}
\end{lstlisting}

\subsection{Szablon, wskaźnik do funkcji czy \texttt{std::function}?}

Jeżeli funkcja korzysta z obiektu wywoływalnego tylko podczas bieżącego
wywołania, często najprościej użyć szablonu. Pozwala to kompilatorowi
znać konkretny typ i daje okazję do \emph{inliningu}. Wskaźnik do funkcji
jest lekki, ale nie przechowa samodzielnie stanu przechwycenia lambdy.
\texttt{std::function} ułatwia natomiast przechowywanie i zastępowanie
heterogenicznych callbacków o wspólnej sygnaturze.

\begin{lstlisting}[caption={Wywołanie generyczne bez type erasure; kompletny program.}]
#include <cassert>
#include <functional>
#include <utility>

template<class F, class... Args>
decltype(auto) call_now(F&& f, Args&&... args) {
    return std::invoke(std::forward<F>(f),
                       std::forward<Args>(args)...);
}

int main() {
    int factor = 3;
    auto f = [factor](int x) { return factor * x; };
    assert(call_now(f, 7) == 21);
}
\end{lstlisting}

\section{Własny mechanizm type erasure: \texttt{Printable}}
\label{sec:cpp14-custom-erasure}

To rozszerzenie pokazuje, co może znajdować się pod gotowym
opakowaniem wymazującym typ. Przy pierwszej lekturze wystarczy
rozumieć, kiedy użyć \texttt{variant}, \texttt{any} lub
\texttt{function}; własną implementację można przeczytać później.

Pokażemy teraz, jak samodzielnie zbudować typ o stałym interfejsie,
który może przechowywać różne wartości drukowalne do strumienia. Celowo
nie wykorzystamy dziedziczenia typów użytkownika po wspólnej klasie
bazowej. Użytkownik powinien móc przekazać \texttt{int},
\texttt{std::string} albo własny typ, dla którego działa
\texttt{operator<<}. Wymagamy przy tym \emph{semantyki wartości}:
kopiowanie \texttt{Printable} ma kopiować przechowywany obiekt.

\paragraph{Idea implementacji.}
Wewnątrz opakowania znajdują się:
\begin{enumerate}
\item abstrakcyjny interfejs \texttt{Concept} z operacjami \texttt{print}
      i \texttt{clone};
\item szablonowy model \texttt{Model<T>}, który przechowuje konkretną
      wartość i realizuje interfejs;
\item wskaźnik \texttt{unique\_ptr<Concept>}, ukrywający konkretny
      typ w miejscu użycia \texttt{Printable}.
\end{enumerate}

\begin{lstlisting}[caption={Własne kopiowalne type erasure (C++20); kompletny program.}]
#include <cassert>
#include <concepts>
#include <memory>
#include <ostream>
#include <sstream>
#include <string>
#include <type_traits>
#include <utility>

class Printable {
    struct Concept {
        virtual ~Concept() = default;
        virtual void print(std::ostream&) const = 0;
        virtual std::unique_ptr<Concept> clone() const = 0;
    };

    template<class T>
    struct Model final : Concept {
        T value;
        template<class U>
        explicit Model(U&& x) : value(std::forward<U>(x)) {}
        void print(std::ostream& os) const override {
            os << value;
        }
        std::unique_ptr<Concept> clone() const override {
            return std::make_unique<Model<T>>(value);
        }
    };

    std::unique_ptr<Concept> self_;

public:
    Printable() = default;

    template<class T>
      requires (!std::same_as<std::remove_cvref_t<T>, Printable> &&
                std::copy_constructible<std::decay_t<T>> &&
                requires(std::ostream& os,
                         const std::decay_t<T>& value) {
                    { os << value } -> std::same_as<std::ostream&>;
                })
    Printable(T&& value)
        : self_(std::make_unique<Model<std::decay_t<T>>>(
              std::forward<T>(value))) {}

    Printable(const Printable& other)
        : self_(other.self_ ? other.self_->clone() : nullptr) {}

    Printable(Printable&&) noexcept = default;

    Printable& operator=(const Printable& other) {
        if (this != &other) {
            Printable tmp(other); // najpierw bezpieczna kopia
            swap(tmp);            // potem nie-rzucajaca zamiana
        }
        return *this;
    }

    Printable& operator=(Printable&&) noexcept = default;

    void swap(Printable& other) noexcept {
        self_.swap(other.self_);
    }

    void print(std::ostream& os) const {
        if (self_) self_->print(os);
        else os << "<empty>";
    }
};

int main() {
    Printable a = 17;
    Printable b = std::string("hello");
    Printable c = b; // gleboka kopia wartosci
    b = 7.5;

    std::ostringstream stream;
    a.print(stream);
    c.print(stream);
    assert(stream.str() == "17hello");
}
\end{lstlisting}

\subsection{Analiza własności i kosztów}

Ten interfejs jest \emph{otwarty}: każdy nowy typ spełniający wymagania
może zostać opakowany bez modyfikowania klasy \texttt{Printable}.
Użytkownik nie musi dziedziczyć po klasie bazowej, co jest ważną zaletą
w porównaniu ze zwykłym polimorfizmem opartym na \texttt{Base}.
\texttt{Model<T>} jest instancjonowany w czasie kompilacji, natomiast
\texttt{print()} jest wywoływane dynamicznie przez ukryty interfejs.
Kopiowanie przez \texttt{clone()} zachowuje rzeczywisty typ wartości.

Wersja edukacyjna wykorzystuje osobną alokację dynamiczną przy
przechowywaniu niepustej wartości i kolejną alokację przy głębokim
kopiowaniu. Nie wykorzystuje \emph{small-buffer optimization}, które
mogłoby przechowywać niewielkie obiekty bezpośrednio wewnątrz opakowania.
W produkcyjnym type erasure trzeba dodatkowo przemyśleć \emph{wyrównanie},
przenoszenie, bezpieczne niszczenie, usuwanie typów niekopiowalnych,
kwalifikatory \texttt{const}, obsługę wyjątków i stabilność ABI.

\paragraph{Gwarancja silna przypisania kopiującego.}
Tworzenie obiektu \texttt{tmp(other)} może zakończyć się wyjątkiem.
Wówczas obiekt docelowy pozostaje niezmieniony. Jeśli kopia się uda,
\texttt{swap(tmp)} jest operacją \texttt{noexcept}; stara zawartość
zostanie zniszczona razem z \texttt{tmp}. Jest to klasyczny idiom
\emph{copy-and-swap}. Alternatywą jest niekopiowalne opakowanie,
które usuwa operacje kopiowania i przechowuje tylko obiekty przenoszalne.

\section{Type erasure a dziedziczenie i polimorfizm statyczny}
\label{sec:cpp14-comparison}

\subsection{Wspólna baza klas}

Gdy typy zostały zaprojektowane jako część wspólnej hierarchii,
\texttt{std::unique\_ptr<Base>} pozwala przechowywać obiekty o różnych
typach dynamicznych i wywoływać ich operacje wirtualne. Implementacje
muszą dziedziczyć po \texttt{Base}, a tworzenie kopii polimorficznej
wymaga własnego protokołu (np. \texttt{clone()}). Z drugiej strony
\texttt{unique\_ptr} jasno wyraża własność i nie kopiuje obiektów
wskutek zwykłego kopiowania uchwytu (ponieważ kopiowanie jest zabronione).

\begin{lstlisting}[caption={Polimorfizm przez bazę i \texttt{unique\_ptr}; kompletny program.}]
#include <cassert>
#include <memory>
#include <string>
#include <vector>

struct Shape {
    virtual ~Shape() = default;
    virtual int area() const = 0;
};
struct Square final : Shape {
    int side;
    explicit Square(int s) : side(s) {}
    int area() const override { return side * side; }
};
struct Rectangle final : Shape {
    int width, height;
    Rectangle(int w, int h) : width(w), height(h) {}
    int area() const override { return width * height; }
};

int main() {
    std::vector<std::unique_ptr<Shape>> shapes;
    shapes.push_back(std::make_unique<Square>(3));
    shapes.push_back(std::make_unique<Rectangle>(2, 5));
    assert(shapes[0]->area() + shapes[1]->area() == 19);
}
\end{lstlisting}

\subsection{Kiedy nie wymazywać typu?}

Jeżeli algorytm potrzebuje konkretnego typu wyłącznie podczas jednego
wywołania, parametr szablonowy z \texttt{concept} jest często prostszy
i pozwala unikać alokacji. Type erasure ma sens przede wszystkim na
granicy interfejsów, gdy wartości o różnych typach muszą być
\emph{przechowywane w jednej strukturze}, przekazywane przez stabilny
interfejs albo zastępowane w trakcie działania programu.

\begin{center}
\small
\begin{tabularx}{\textwidth}{@{}lXX@{}}
\toprule
\textbf{Rozwiązanie} & \textbf{Zalety} & \textbf{Ograniczenia} \\
\midrule
\texttt{variant<Ts...>} & Typy znane statycznie, sprawdzanie przypadków, lokalne przechowywanie wartości & Zamknięty katalog typów, złożoność odwiedzania i rozmiar największej alternatywy \\
\texttt{any} & Dowolna kopiowalna wartość, wspólny typ pojemnika & Trzeba odzyskać dokładny typ; brak wspólnego interfejsu operacji \\
\texttt{function} & Jedna sygnatura dla różnych callbacków & Kopiowalność (C++20), pośrednictwo, możliwe alokacje \\
Własne type erasure & Otwarty katalog modeli bez wspólnej bazy typów klienta & Złożona implementacja, zarządzanie pamięcią i wyjątkami \\
\texttt{unique\_ptr<Base>} & Jawna hierarchia, dziedziczone operacje, kontrola własności & Wymaga wspólnej bazy; operacje wirtualne \\
Szablon z \texttt{concept} & Typ znany kompilatorowi, możliwość optymalizacji & Wiele instancjacji; zwykle nie przechowa różnych typów w jednej kolekcji \\
\bottomrule
\end{tabularx}
\end{center}

\section{Wydajność: układ pamięci i koszty dynamicznej dyspozycji}
\label{sec:cpp14-performance}

Nie ma uniwersalnej odpowiedzi na pytanie, który mechanizm jest
najszybszy. Poniżej wymieniamy dominujące koszty; faktyczne zachowanie
zależy od rozmiaru typów, optymalizacji kompilatora, struktury programu
i rozkładu typów w danych.

\begin{itemize}
\item \texttt{optional<T>} dodaje znacznik obecności; dostęp do zawartości
      może wymagać sprawdzenia warunku.
\item \texttt{variant<Ts...>} przechowuje wystarczającą ilość miejsca
      na największy typ i indeks aktywnej alternatywy; odwiedzający
      może być implementowany przez instrukcję wyboru lub tablicę
      funkcji, a szczegóły zależą od kompilatora.
\item \texttt{any} przechowuje metadane typu oraz wartość, potencjalnie
      w osobno alokowanym buforze.
\item \texttt{function} może wymagać pośredniego wywołania i przechowania
      stanu lambdy; czasami da się uniknąć alokacji.
\item Własna implementacja z \texttt{unique\_ptr<Concept>} wykonuje
      alokację dynamiczną i pośrednie wywołanie wirtualne.
\end{itemize}

Przede wszystkim trzeba mierzyć \emph{rzeczywisty} przypadek użycia.
Jednorazowy koszt konstruowania callbacku może być nieistotny, gdy jego
wykonanie trwa długo, ale może dominować w pętli składającej się z
milionów krótkich operacji.

\subsection{Test bez mylenia czasu kompilacji z wykonaniem}

Rzetelne porównanie wymaga rozgrzania kodu, wykorzystania odpowiedniego
zestawu danych, uwzględnienia optymalizacji i sprawdzenia, czy kompilator
nie usunął mierzonej pracy. Warto rozdzielać koszt \emph{tworzenia}
opakowania od kosztu \emph{wywołania} i analizować nie tylko średni
czas, ale też rozkład wyników oraz liczbę alokacji. Benchmarku
porównującego dwa mechanizmy nie traktujemy jako dowodu ich przewagi
we wszystkich zastosowaniach.

\section{Granice bezpieczeństwa: wyjątki, czas życia i współbieżność}

\paragraph{Wyjątki.}
\texttt{optional::value()}, \texttt{std::get} dla wariantu, \texttt{any\_cast}
i wywołanie pustego \texttt{std::function} mogą zgłaszać odpowiednie
wyjątki. Konstrukcja zawartości i przypisanie także mogą rzucać,
ponieważ wykonują operacje typu przechowywanego. Należy rozróżniać
\emph{brak wartości} (normalny stan \texttt{optional}) od
\emph{niepowodzenia operacji}, któremu warto nadać osobną reprezentację.

\paragraph{Czas życia.}
\texttt{variant}, \texttt{optional} i \texttt{any} posiadają własne
przechowywane wartości. Można jednak umieścić w nich wskaźniki,
\texttt{reference\_wrapper} albo inne nieposiadające własności uchwyty;
wtedy odpowiedzialność za czas życia wskazywanych danych pozostaje po
stronie programisty. To samo dotyczy referencji przechwyconych przez
lambdę umieszczoną w \texttt{std::function}.

\paragraph{Wątki.}
Żaden z opisanych typów nie czyni dowolnych współbieżnych modyfikacji
tego samego obiektu bezpiecznymi. Przykładowo równoczesny zapis do
jednego \texttt{variant} i odczyt jego aktywnej zawartości bez
synchronizacji mogą powodować wyścig danych. Gdy callback jest
uruchamiany równolegle, trzeba osobno określić zasady synchronizacji
stanu przechwyconego przez lambdę.

\section{Nowości po C++20 i świadomy wybór standardu}

W całym rozdziale przykłady są zgodne z C++20. W nowszych standardach
pojawiają się udogodnienia, które warto znać, ale których \emph{nie wolno}
zakładać przy kompilacji z \texttt{-std=c++20}:
\begin{itemize}
\item \textbf{C++23:} \texttt{std::expected<T,E>} --- wynik albo
      jawna wartość błędu;
\item \textbf{C++23:} \texttt{std::move\_only\_function} --- opakowanie
      dla obiektów wywoływalnych, które nie muszą być kopiowalne;
\item \textbf{C++23:} operacje monadyczne dla \texttt{std::optional}:
      \texttt{and\_then}, \texttt{transform} i \texttt{or\_else};
\item \textbf{C++26:} \texttt{std::function\_ref} --- niewłaścicielski
      widok na obiekt wywoływalny, którego czas życia musi zostać
      zapewniony przez wywołującego.
\end{itemize}

Dostępność funkcji bibliotecznych w konkretnym kompilatorze i jego
bibliotece standardowej może być opóźniona w stosunku do daty
opublikowania standardu.

\section{Zasady projektowania i typowe błędy}
\label{sec:cpp14-guidelines}

\begin{enumerate}
\item Wybieraj \texttt{optional<T>}, gdy brak wyniku jest normalny i
      nie potrzebujesz szczegółowego typu błędu.
\item Wybieraj \texttt{variant}, gdy zestaw przypadków jest znany
      i zamknięty; obsługuj je przez \texttt{visit}, jeśli pozwala to
      uniknąć łańcucha ręcznych kontroli typu.
\item Stosuj \texttt{any} tylko tam, gdzie świadomie rezygnujesz
      ze statycznie określonego zestawu typów i masz jasny protokół
      odzyskiwania ich wartości.
\item Używaj \texttt{std::function} do \emph{przechowywanych},
      kopiowalnych callbacków; dla funkcji wykonywanej natychmiast
      rozważ parametr szablonowy.
\item Nie zakładaj, że \texttt{std::function} zawsze alokuje ani
      że \texttt{std::variant} zawsze jest szybszy. Zmierz przypadek
      zastosowania.
\item Nie pozostawiaj referencji lub wskaźników do obiektów, których
      czas życia zakończy się przed użyciem przechowywanego callbacku.
\item Unikaj \texttt{std::get} i \texttt{any\_cast} w sytuacjach,
      w których nie masz gwarancji typu; do sprawdzania wykorzystuj
      \texttt{get\_if} i wskaźnikowy \texttt{any\_cast}.
\item Nie używaj \texttt{std::any} do bezpośredniego opakowania
      typów przenoszalnych, ale niekopiowalnych, np. \texttt{unique\_ptr}.
\item Rozdziel decyzję o \emph{otwartym zbiorze implementacji} od
      decyzji o \emph{dynamicznym czasie życia i alokacji}.
\end{enumerate}

\section*{Podsumowanie rozdziału}
\addcontentsline{toc}{section}{Podsumowanie rozdziału 10}

\texttt{optional} modeluje brak lub obecność wartości, a
\texttt{variant} modeluje zamkniętą sumę przypadków. \texttt{any}
ukrywa typ kopiowalnej wartości, nie narzucając jej operacji.
\texttt{function} ukrywa typ obiektu wywoływalnego przy zachowaniu
ustalonej sygnatury. Własne type erasure pozwala opakować otwartą
rodzinę obcych typów bez wymagania wspólnej klasy bazowej, ale wymaga
starannego określenia semantyki własności, kopiowania i wyjątków.
Dobry projekt zaczyna się więc od określenia kontraktu: \emph{które
przypadki istnieją, kto posiada obiekty, jakie operacje udostępniamy
oraz czy zbiór typów ma pozostać otwarty}.

\section*{Dokumentacja i literatura}
\addcontentsline{toc}{section}{Dokumentacja i literatura rozdziału 10}
\begin{itemize}
\item \url{https://en.cppreference.com/w/cpp/utility/optional} --- \texttt{std::optional}.
\item \url{https://en.cppreference.com/w/cpp/utility/variant} --- \texttt{std::variant}.
\item \url{https://en.cppreference.com/w/cpp/utility/any} --- \texttt{std::any}.
\item \href{https://en.cppreference.com/w/cpp/utility/functional/function}{cppreference: \texttt{std::function}} --- dokumentacja opakowania wywołań.
\item Bjarne Stroustrup, \emph{A Tour of C++}, wyd. 3, Addison--Wesley, 2022.
\item Klaus Iglberger, \emph{C++ Software Design}, O'Reilly, 2022.
\end{itemize}
